\section{Θέματα Α}
%=========== 1ο ΘΕΜΑ Α ===========
\begin{thema}{A}
\begin{erwthma}
    \item Δίνεται μια συνάρτηση $ f $ ορισμένη σε ένα διάστημα $ \varDelta $ και έστω $ F $ μια παράγουσα της $ f $ στο $ \varDelta $. Να αποδείξετε ότι \begin{alist} \item όλες οι συναρτήσεις της μορφής \[ G(x)=F(x)+c \] είναι παράγουσες της $ f $ στο $ \varDelta $ και ότι \item κάθε άλλη παράγουσα $ G $ της $ f $ στο $ \varDelta $ παίρνει τη μορφή \[ G(x)=F(x)+c \] για κάθε $ x\in\varDelta $, με $ c\in\mathbb{R} $. \end{alist}
    \item Τι ονομάζουμε πραγματική συνάρτηση με πεδίο ορισμού το σύνολο $A$?
    \item Πότε μία συνάρτηση $f$ ονομάζεται γνησίως αύξουσα σε ένα διάστημα του πεδίου ορισμού της?
    \item Να χαρακτηρίσετε τις προτάσεις που ακολουθούν, γράφοντας στο τετράδιό σας, δίπλα στο γράμμα που αντιστοιχεί σε κάθε πρόταση, τη λέξη \textbf{Σωστό}, αν η πρόταση είναι σωστή, ή \textbf{Λάθος}, αν η πρόταση είναι λανθασμένη.
    \begin{alist}
        \item Το σύνολο τιμών μιας συνεχούς και μη σταθερής συνάρτησης $f$ σε κλειστό διάστημα είναι πάντα κλειστό διάστημα.
        \item Αν μια παραγωγίσιμη συνάρτηση δεν είναι γνησίως αύξουσα, τότε η παράγωγός της δεν μπορεί να είναι παντού θετική.
        \item Κάθε γνησίως μονότονη συνάρτηση είναι $1-1$.
        \item Κάθε συνεχής συνάρτηση σε ένα σημείο $x_0$, είναι και παραγωγίσιμη σε αυτό.
        \item Το Θ.Μ.Τ. προϋποθέτει ότι η συνάρτηση είναι παραγωγίσιμη και στα άκρα του διαστήματος.
    \end{alist}
    \item Αν $x_0$ είναι εσωτερικό σημείο και τοπικό ακρότατο της $f$, και η $f$ παραγωγίζεται στο $x_0$, τότε:
    \begin{tasks}[label=\let\textdexiakeraia\relax\alph*.](4)
        \task $f'(x_0) > 0$
        \task $f'(x_0) < 0$
        \task $f'(x_0) = 0$
        \task Δεν εξάγεται συμπέρασμα
    \end{tasks}
\end{erwthma}
\end{thema}
%=========== ΤΕΛΟΣ - 1ο ΘΕΜΑ Α ===========


%=========== 2ο ΘΕΜΑ Α ===========
\begin{thema}{A}
\begin{erwthma}
    \item Να αποδείξετε ότι αν μια συνάρτηση $f$ είναι παραγωγίσιμη σ’ ένα σημείο {$ x_0 $} τότε είναι και συνεχής στο σημείο αυτό.
    \item Πότε μια συνάρτηση $f:A\to\mathbb{R}$ ονομάζεται συνάρτηση $1-1$?
    \item Πότε μια συνάρτηση $f$ με πεδίο ορισμού $A$ παρουσιάζει τοπικό μέγιστο σε ένα σημείο $x_0\in A$?
    \item Να χαρακτηρίσετε τις προτάσεις που ακολουθούν, γράφοντας στο τετράδιό σας, δίπλα στο γράμμα που αντιστοιχεί σε κάθε πρόταση, τη λέξη \textbf{Σωστό}, αν η πρόταση είναι σωστή, ή \textbf{Λάθος}, αν η πρόταση είναι λανθασμένη.
    \begin{alist}
        \item Η παράγωγος $\left(\dint_a^\beta f(x) \d x \right)'$ ισούται με $ f(x)$.
        \item Το ολοκλήρωμα του γινομένου ισούται πάντοτε με το γινόμενο των ολοκληρωμάτων.
        \item Κάθε ολικό μέγιστο είναι σίγουρα μεγαλύτερο από οποιοδήποτε τοπικό ελάχιστο.
        \item Αν $\dint_a^\beta f(x) \d x \ge 0$, τότε υποχρεωτικά η $f$ είναι μη αρνητική παντού στο $[a, \beta]$.
        \item Η έννοια της συνέχειας στο $x_0$ προϋποθέτει ότι το $x_0$ ανήκει στο πεδίου ορισμού.
    \end{alist}
    \item Κατά την ολοκλήρωση με αντικατάσταση $u = g(x)$, αλλάζουν:
    \begin{tasks}[label=\let\textdexiakeraia\relax\alph*.](2)
        \task Μόνο η μεταβλητή ολοκλήρωσης.
        \task H μεταβλητή ολοκλήρωσης και τα όρια ολοκλήρωσης.
        \task Μόνο τα όρια ολοκλήρωσης.
        \task Τίποτα.
    \end{tasks}
\end{erwthma}
\end{thema}
%=========== ΤΕΛΟΣ - 2ο ΘΕΜΑ Α ===========


%=========== 3ο ΘΕΜΑ Α ===========
\begin{thema}{A}
\begin{erwthma}
    \item Αν οι συναρτήσεις $ f,g $ είναι παραγωγίσιμες στο $ x_0 $, τότε η συνάρτηση $ f+g $ είναι παραγωγίσιμη στο $ x_0 $ και ισχύει: \[ (f+g)'(x_0)=f'(x_0)+g'(x_0) \]
    \item Έστω μια συνάρτηση $f$ ορισμένη σε ένα σύνολο $A$. Τι ονομάζεται πρώτη παράγωγος της $f$?
    \item Πότε μια συνάρτηση $f$ ονομάζεται συνεχής? Πότε συνεχής σε κλειστό διάστημα $[a, \beta]$?
    \item Να χαρακτηρίσετε τις προτάσεις που ακολουθούν, γράφοντας στο τετράδιό σας, δίπλα στο γράμμα που αντιστοιχεί σε κάθε πρόταση, τη λέξη \textbf{Σωστό}, αν η πρόταση είναι σωστή, ή \textbf{Λάθος}, αν η πρόταση είναι λανθασμένη.
    \begin{alist}
        \item Η παράγωγος σταθερής συνάρτησης είναι $0$.
        \item Ισχύει $\dint_a^\beta f(t) dt + \dint_\beta^a f(x) \d x = 0$.
        \item Η γραφική παράσταση μιας συνάρτησης μπορεί να τέμνει την κατακόρυφη ασύμπτωτή της.
        \item Αν $\lim\limits_{x\to x_0}{f(x)}=l$ και $\lim\limits_{x\to x_0}{g(x)}=m,\ l,m\in\mathbb{R}$ και $f(x)<g(x)$ κοντά στο $x_0$ τότε θα ισχύει $l\leq m$.
        \item Αν $f$ είναι $1-1$, η παράγωγος της δεν έχει ρίζες.
    \end{alist}
    \item Η παράγωγος της σύνθετης συνάρτησης $f(g(x))$ ισούται με:
    \begin{tasks}[label=\let\textdexiakeraia\relax\alph*.](2)
        \task $f'(x) \cdot g'(x)$
        \task $f'(g(x)) \cdot g'(x)$
        \task $f(g'(x))$
        \task $f'(g(x)) + g'(x)$
    \end{tasks}
\end{erwthma}
\end{thema}
%=========== ΤΕΛΟΣ - 3ο ΘΕΜΑ Α ===========


%=========== 4ο ΘΕΜΑ Α ===========
\begin{thema}{A}
\begin{erwthma}
    \item Να αποδείξετε ότι η συνάρτηση $ f(x)=\syf{x} $ είναι παραγωγίσιμη στο σύνολο $ A=\{x\in\mathbb{R}|\hm{x}\neq 0\} $ και ισχύει $ f'(x)=-\dfrac{1}{\hm^2{x}} $.
    \item Τι ονομάζεται γραφική παράσταση μιας συνάρτησης $f$?
    \item Πότε μια συνάρτηση $f$, με πεδίο ορισμού $A$, λέγεται παραγωγίσιμη?
    \item Να χαρακτηρίσετε τις προτάσεις που ακολουθούν, γράφοντας στο τετράδιό σας, δίπλα στο γράμμα που αντιστοιχεί σε κάθε πρόταση, τη λέξη \textbf{Σωστό}, αν η πρόταση είναι σωστή, ή \textbf{Λάθος}, αν η πρόταση είναι λανθασμένη.
    \begin{alist}
        \item Αν $f$ συνεχή στο κλειστό $[a, \beta]$ και δεν έχει κρίσιμα σημεία, τα ακρότατά της εντοπίζονται υποχρεωτικά στα άκρα.
        \item Αν $\lim\limits_{x\to x_0} |f(x)| = 1$, τότε $\lim\limits_{x\to x_0} f(x) = 1$ ή $\lim\limits_{x\to x_0} f(x) = -1$.
        \item Αν ισχύει $f''(x) > 0$ για κάθε $\in\varDelta$, τότε υποχρεωτικά η $f(x)$ είναι γνησίως αύξουσα στο διάστημα $\varDelta$.
        \item Αν η $f$ είναι συνεχής με $\dint_a^\beta f(x) \d x = 0$ και $f(x) \ge 0$, τότε $f(x) = 0$ για κάθε $x \in [a, \beta]$.
        \item Αν $f$ άρτια, τότε $\dint_{-a}^a f(x) \d x = 2 \dint_0^a f(x) \d x$.
    \end{alist}
    \item Ο τύπος ολοκλήρωσης κατά παράγοντες είναι:
    \begin{tasks}[label=\let\textdexiakeraia\relax\alph*.](1)
        \task $\dint_a^\beta f(x)g'(x) \d x = f(\beta)g(\beta) - \dint_a^\beta f'(x)g(x) \d x$
        \task $\dint_a^\beta f(x)g'(x) \d x = [f(x)g(x)]_a^\beta - \dint_a^\beta f'(x)g(x) \d x$
        \task $\dint_a^\beta f'(x)g'(x) \d x = f(x)g(x)$
        \task $\dint_a^\beta (f(x)g(x))' \d x = f(x) g(x)$
    \end{tasks}
\end{erwthma}
\end{thema}
%=========== ΤΕΛΟΣ - 4ο ΘΕΜΑ Α ===========


%=========== 5ο ΘΕΜΑ Α ===========
\begin{thema}{A}
\begin{erwthma}
    \item Δίνεται μια συνάρτηση $ f $ η οποία είναι παραγωγίσιμη σ’ ένα διάστημα $ (a,\beta) $, με εξαίρεση ίσως ένα σημείο του $ x_0 $, στο οποίο όμως είναι συνεχής. Να αποδείξετε ότι αν $ f'(x)>0 $ για κάθε $ x\in(a,x_0) $ και $ f'(x)<0 $ για κάθε $ x\in(x_0,\beta) $ τότε η $ f $ παρουσιάζει τοπικό μέγιστο στο $ x_0 $. 
    \item Πότε δύο συναρτήσεις $f, g$ ονομάζονται ίσες?
    \item Ποιες είναι οι πιθανές θέσεις σημείων καμπής μια παραγωγίσιμης συνάρτησης $f$?
    \item Να χαρακτηρίσετε τις προτάσεις που ακολουθούν, γράφοντας στο τετράδιό σας, δίπλα στο γράμμα που αντιστοιχεί σε κάθε πρόταση, τη λέξη \textbf{Σωστό}, αν η πρόταση είναι σωστή, ή \textbf{Λάθος}, αν η πρόταση είναι λανθασμένη.
    \begin{alist}
        \item Αν $f$ $1-1$ και παραγωγίσιμη, η αντίστροφη είναι επίσης παραγωγίσιμη στο πεδίο αντίστοιχων τιμών, αρκεί $f'(x) \ne 0$.
        \item Η συνάρτηση $f(x)=x^3$ έχει εξίσωση εφαπτομένης στο $x=0$ τον άξονα $x'x$.
        \item H συνάρτηση $f(x)=x^2$ δεν είναι $1-1$ στο $\mathbb{R}$.
        \item Αν μια συνάρτηση είναι κυρτή σένα διάστημα $\Delta$ του πεδίου ορισμού της, η εφαπτομένη της βρίσκεται κάτω από τη γραφική παράσταση, με εξαίρεση το σημείο επαφής.
        \item Αν $\lim\limits_{x \to +\infty} \dfrac{f(x)}{x} = \lambda$ και $\lim\limits_{x \to +\infty} (f(x)-\lambda x) = \mu$, τότε η $y=\lambda x + \mu$ είναι πλάγια (ή οριζόντια) ασύμπτωτη στο $+\infty$.
    \end{alist}
    \item Η ιδιότητα $\dint_a^\beta f(x) \d x = - \dint_\beta^a f(x) \d x$ ισχύει:
    \begin{tasks}[label=\let\textdexiakeraia\relax\alph*.](2)
        \task Μόνο αν $a > \beta$.
        \task Μόνο αν $f(x) \ge 0$.
        \task Εφόσον η $f$ είναι συνεχής, για οποιαδήποτε $a, \beta$.
        \task Ποτέ.
    \end{tasks}
\end{erwthma}
\end{thema}
%=========== ΤΕΛΟΣ - 5ο ΘΕΜΑ Α ===========


%=========== 6ο ΘΕΜΑ Α ===========
\begin{thema}{A}
\begin{erwthma}
    \item Να διατυπώσετε το θεώρημα \eng{Bolzano} και να δώσετε τη γεωμετρική ερμηνεία του.
    \item Πώς ορίζονται οι πράξεις της πρόσθεσης $f+g$ και του πολλαπλασιασμού $f \cdot g$ δύο συναρτήσεων? Ποιο το πεδίο ορισμού τους?
    \item Πότε λέμε ότι μια συνάρτηση $f$ στρέφει τα κοίλα προς τα άνω (ή είναι κυρτή) στο $\Delta$?
    \item Να χαρακτηρίσετε τις προτάσεις που ακολουθούν, γράφοντας στο τετράδιό σας, δίπλα στο γράμμα που αντιστοιχεί σε κάθε πρόταση, τη λέξη \textbf{Σωστό}, αν η πρόταση είναι σωστή, ή \textbf{Λάθος}, αν η πρόταση είναι λανθασμένη.
    \begin{alist}
        \item Η $f(x) = \ln x$ έχει παράγουσα την $F(x) = x\ln x - x + c$.
        \item Μια συνάρτηση παραγωγίσιμη στο $x_0\in D_f$ είναι πάντα συνεχής στο $x_0\in D_f$.
        \item Αν $\dint_a^\beta f(x) \d x < 0$, τότε δεν υπάρχει κανένα $x\in[a, \beta]$ ώστε $f(x)\ge 0$.
        \item Ο τύπος ολοκλήρωσης κατά παράγοντες είναι $\dint_a^\beta f(x)g'(x) \d x = [f(x)g(x)]_a^\beta - \dint_a^\beta f'(x)g'(x) \d x$.
        \item Αν η $f''(x)$ αλλάζει πρόσημο εκατέρωθεν του $x_0$ (όπου ορίζεται η εφαπτομένη), το $x_0$ είναι σημείο καμπής.
    \end{alist}
    \item Η κατακόρυφη ασύμπτωτη εμφανίζεται στα σημεία $x_0$ όπου:
    \begin{tasks}[label=\let\textdexiakeraia\relax\alph*.](2)
        \task Ο αριθμητής μηδενίζεται και ο παρονομαστής όχι.
        \task Η συνάρτηση παρουσιάζει ελάχιστο.
        \task Κάποιο πλευρικό όριο καταλήγει σε $+\infty$ ή $-\infty$.
        \task Η συνάρτηση είναι παραγωγίσιμη.
    \end{tasks}
\end{erwthma}
\end{thema}
%=========== ΤΕΛΟΣ - 6ο ΘΕΜΑ Α ===========


%=========== 7ο ΘΕΜΑ Α ===========
\begin{thema}{A}
\begin{erwthma}
    \item Δίνονται δύο συναρτήσεις $ f,g $ ορισμένες σε ένα διάστημα $ \varDelta $. Να αποδείξετε ότι αν \begin{itemize} \item οι συναρτήσεις $ f,g $ είναι συνεχείς στο διάστημα $ \varDelta $ και \item $ f'(x)=g'(x) $ σε κάθε \emph{εσωτερικό} σημείο του $ \varDelta $ \end{itemize} τότε υπάρχει σταθερά $ c $ τέτοια ώστε να ισχύει $ f(x)=g(x)+c $ για κάθε $ x\in\varDelta $.
    \item Έστω οι συναρτήσεις $f$ και $g$. Πώς ορίζεται η σύνθεση $g \circ f$ και ποιο είναι το πεδίο ορισμού της?
    \item Πότε η ευθεία $x=x_0$ ονομάζεται κατακόρυφη ασύμπτωτη της γραφικής παράστασης μιας συνάρτησης $f$?
    \item Να χαρακτηρίσετε τις προτάσεις που ακολουθούν, γράφοντας στο τετράδιό σας, δίπλα στο γράμμα που αντιστοιχεί σε κάθε πρόταση, τη λέξη \textbf{Σωστό}, αν η πρόταση είναι σωστή, ή \textbf{Λάθος}, αν η πρόταση είναι λανθασμένη.
    \begin{alist}
        \item Το ολοκλήρωμα $\dint_a^\beta |f(x)| \d x$ αντιπροσωπεύει πάντοτε το εμβαδόν του χωρίου μεταξύ της $C_{f}$, του άξονα $x'x$ και των $x=a$ και $x=\beta$, ακόμα και αν δεν γνωρίζουμε ποιο είναι μεγαλύτερο ($a$ ή $\beta$).
        \item Κάθε συνεχής συνάρτηση έχει σταθερό πρόσημο ανάμεσα σε δύο διαδοχικές ρίζες της.
        \item Κάθε περιοδική μη σταθερή συνάρτηση δεν έχει όριο στο άπειρο.
        \item Αν η $f$ είναι συνεχής και $f(x) \ge 0$ για κάθε $x \in [a, \beta]$ με $\dint_a^\beta f(x) \d x = 0$, τότε υποχρεωτικά $f(x) = 0$ σε όλο το διάστημα.
        \item Οι γραφικές παραστάσεις $C_f, C_{f^{-1}}$ είναι συμμετρικές ως προς τη διχοτόμο $y=x$.
    \end{alist}
    \item Ο γεωμετρικός ορισμός της $f'(x_0)$ ταυτίζεται με:
    \begin{tasks}[label=\let\textdexiakeraia\relax\alph*.](2)
        \task Τον άξονα συμμετρίας της συνάρτησης.
        \task Το εμβαδόν κάτω από την γραφική παράσταση.
        \task Το συντελεστή διεύθυνσης $\lambda$ της εφαπτομένης της $C_f$ στο $x_0$.
        \task Την ελάχιστη τιμή της $f$.
    \end{tasks}
\end{erwthma}
\end{thema}
%=========== ΤΕΛΟΣ - 7ο ΘΕΜΑ Α ===========


%=========== 8ο ΘΕΜΑ Α ===========
\begin{thema}{A}
\begin{erwthma}
    \item Έστω η σταθερή συνάρτηση $ f(x)=c, c\in\mathbb{R} $. Να αποδείξετε ότι η συνάρτηση $ f $ είναι παραγωγίσιμη στο $ \mathbb{R} $ και ισχύει $ f'(x)=0 $, δηλαδή $ (c)'=0 $.
    \item Πότε μια συνάρτηση λέγεται γνησίως φθίνουσα σε ένα διάστημα $\Delta$?
    \item Πότε η ευθεία $y=y_0$ ονομάζεται οριζόντια ασύμπτωτη της γραφικής παράστασης μιας συνάρτησης $f$?
    \item Να χαρακτηρίσετε τις προτάσεις που ακολουθούν, γράφοντας στο τετράδιό σας, δίπλα στο γράμμα που αντιστοιχεί σε κάθε πρόταση, τη λέξη \textbf{Σωστό}, αν η πρόταση είναι σωστή, ή \textbf{Λάθος}, αν η πρόταση είναι λανθασμένη.
    \begin{alist}
        \item Αν η συνάρτηση $f$ έχει τοπικό μέγιστο στο σημείο $x_0$, τότε υποχρεωτικά ισχύει $f'(x_0) = 0$.
        \item Αν $x_0$ είναι ακρότατο σε κλειστό άκρο διαστήματος, τότε ισχύει το \eng{Fermat}.
        \item Αν $\lim\limits_{x\to x_0} f(x) = -\infty$, τότε $\lim\limits_{x\to x_0} |f(x)| = +\infty$.
        \item Αν $f(x) < g(x)$ για κάθε $x \in [a, \beta]$ με $f,g$ συνεχείς, τότε $\dint_a^\beta f(x) \d x < \dint_a^\beta g(x) \d x$.
        \item Αν $f'(x) = 0$ σε ένωση ξένων διαστημάτων, τότε η $f$ μπορεί να παίρνει διαφορετικές σταθερές τιμές σε κάθε διάστημα.
    \end{alist}
    \item Αν $f$ άρτια και ολοκληρώσιμη στο $[-a, a]$ (με $a > 0$), τότε:
    \begin{tasks}[label=\let\textdexiakeraia\relax\alph*.](2)
        \task $\dint_{-a}^{a} f(x) \d x = 0$
        \task $\dint_{-a}^{a} f(x) \d x = 2\dint_0^{a} f(x) \d x$
        \task $\dint_{-a}^{a} f(x) \d x = a \cdot f(0)$
        \task H $f$ δεν είναι ολοκληρώσιμη.
    \end{tasks}
\end{erwthma}
\end{thema}
%=========== ΤΕΛΟΣ - 8ο ΘΕΜΑ Α ===========


%=========== 9ο ΘΕΜΑ Α ===========
\begin{thema}{A}
\begin{erwthma}
    \item Να διατυπώσετε και να δώσετε τη γεωμετρική ερμηνεία του θεωρήματος \eng{Rolle}.
    \item Τι ονομάζουμε ολικό μέγιστο και τι ολικό ελάχιστο μιας συνάρτησης $f$?
    \item Πότε η ευθεία $y=\lambda x+\beta$ ονομάζεται πλάγια ασύμπτωτη της γραφικής παράστασης μιας συνάρτησης $f$ στο $+\infty$?
    \item Να χαρακτηρίσετε τις προτάσεις που ακολουθούν, γράφοντας στο τετράδιό σας, δίπλα στο γράμμα που αντιστοιχεί σε κάθε πρόταση, τη λέξη \textbf{Σωστό}, αν η πρόταση είναι σωστή, ή \textbf{Λάθος}, αν η πρόταση είναι λανθασμένη.
    \begin{alist}
        \item Η παράγωγος πολυωνύμου 3ου βαθμού έχει ασύμπτωτες.
        \item Αν μια συνάρτηση δεν είναι $1-1$, τότε δεν είναι γνησίως μονότονη.
        \item Αν η $f'$ είναι συνεχής, τότε $\dint_a^\beta f'(x) \d x = f(\beta) - f(a)$.
        \item Μια συνάρτηση μπορεί να έχει άπειρες κατακόρυφες ασύμπτωτες.
        \item Αν μία συνάρτηση $f$ είναι συνεχής στο $[a, \beta]$ και δεν έχει ρίζα, τότε $f(a)f(\beta) \ge 0$.
    \end{alist}
    \item Αν η συνάρτηση είναι παραγωγίσιμη σε ένα σημείο $x_0$, τότε:
    \begin{tasks}[label=\let\textdexiakeraia\relax\alph*.](2)
        \task Υποχρεωτικά $f'(x_0) = 0$.
        \task Είναι οπωσδήποτε συνεχής στο $x_0$.
        \task Μπορεί να έχει διαφορετικά πλευρικά όρια.
        \task Παρουσιάζει ακρότατο.
    \end{tasks}
\end{erwthma}
\end{thema}
%=========== ΤΕΛΟΣ - 9ο ΘΕΜΑ Α ===========


%=========== 10ο ΘΕΜΑ Α ===========
\begin{thema}{A}
\begin{erwthma}
    \item Θεωρούμε μια συνάρτηση $f$ ορισμένη σε ένα κλειστό διάστημα $[a,\beta]$. Αν
\begin{itemize}
\item η $f$ συνεχής στο κλειστό διάστημα $[a,\beta]$ και 
\item $f(a)\neq f(\beta)$
\end{itemize}
Να αποδείξετε ότι υπάρχει τουλάχιστον ένα $ x_0\in(a,\beta) $ ώστε για κάθε αριθμό $ \eta $ μεταξύ των $ f(a),f(\beta) $ να ισχύει $ f(x_0)=\eta $.
    \item Έστω μια συνάρτηση $f$ με πεδίο ορισμού $A$. Τι ονομάζεται δεύτερη παράγωγος της $f$ και πώς ορίζεται η $\nu$-οστή παράγωγος της $f$?
    \item Τι ονομάζουμε αρχική συνάρτηση ή παράγουσα μιας συνάρτησης $f$ σε ένα διάστημα $\Delta$?
    \item Να χαρακτηρίσετε τις προτάσεις που ακολουθούν, γράφοντας στο τετράδιό σας, δίπλα στο γράμμα που αντιστοιχεί σε κάθε πρόταση, τη λέξη \textbf{Σωστό}, αν η πρόταση είναι σωστή, ή \textbf{Λάθος}, αν η πρόταση είναι λανθασμένη.
    \begin{alist}
        \item Αν $\dint_a^\beta f(x) \d x = 0$ και η $f$ παίρνει και θετικές και αρνητικές τιμές, το εμβαδόν του χωρίου που περικλείεται από την $C_f$, τον άξονα $x'x$ και τις ευθείες $x=a$ και $x=\beta$ είναι θετικό.
        \item Αν υπάρχει τουλάχιστον ένα $\xi \in (a, \beta)$ με $f'(\xi)=0$, τότε υποχρεωτικά $f(a)=f(\beta)$.
        \item Η συνάρτηση $f(x) = \sqrt{x}$ είναι παραγωγίσιμη στο $x=0$.
        \item Το πεδίο ορισμού του αθροίσματος $f+g$ είναι πάντοτε το $D_f\cap D_g$.
        \item Αν $\dint_a^\beta f(x) \d x = \dint_a^\beta g(x) \d x$, τότε υποχρεωτικά συνεπάγεται ότι $f(x) = g(x)$ σε όλο το $[a, \beta]$.
    \end{alist}
    \item Αν η παράγωγος $f'(x_0)$ δεν υπάρχει, ενώ η $f$ είναι συνεχής στο $x_0$, τότε στο $x_0$:
    \begin{tasks}[label=\let\textdexiakeraia\relax\alph*.](2)
        \task Η $f$ δεν μπορεί να έχει ακρότατο.
        \task Η $f$ μπορεί να έχει ακρότατο (π.χ. $f(x)=|x|$ στο $0$).
        \task Η $f$ είναι αναγκαστικά σταθερή κοντά στο $x_0$.
        \task Η $f$ δεν ορίζεται στο $x_0$.
    \end{tasks}
\end{erwthma}
\end{thema}
%=========== ΤΕΛΟΣ - 10ο ΘΕΜΑ Α ===========


%=========== 11ο ΘΕΜΑ Α ===========
\begin{thema}{A}
\begin{erwthma}
    \item Να αποδείξετε ότι η συνάρτηση $ f(x)=x^a $ με $ a\in\mathbb{R}-\mathbb{Z} $ είναι παραγωγίσιμη στο $ (0,+\infty) $ με $ f'(x)=ax^{a-1} $.
    \item Τι ονομάζουμε ρυθμό μεταβολής ενός ποσού $y=f(x)$ ως προς ένα ποσό $x$ σε ένα σημείο $x_0$?
    \item Πότε μια συνάρτηση $f$ με πεδίο ορισμού $A$ παρουσιάζει τοπικό ελάχιστο σε ένα σημείο $x_0\in A$?
    \item Να χαρακτηρίσετε τις προτάσεις που ακολουθούν, γράφοντας στο τετράδιό σας, δίπλα στο γράμμα που αντιστοιχεί σε κάθε πρόταση, τη λέξη \textbf{Σωστό}, αν η πρόταση είναι σωστή, ή \textbf{Λάθος}, αν η πρόταση είναι λανθασμένη.
    \begin{alist}
        \item Αν το $x_0$ είναι σημείο καμπής και υπάρχει η δεύτερη παράγωγος, τότε υποχρεωτικά $f''(x_0) = 0$.
        \item Μια συνάρτηση $f$ η οποία είναι άρτια και έχει πεδίο ορισμού το $\mathbb{R}$, δεν μπορεί να είναι $1-1$.
        \item Το πεδίο ορισμού του πηλίκου $\dfrac{f}{g}$ δύο συναρτήσεων $f,g$ είναι το $D_f\cap D_g$.
        \item Ένα τοπικό μέγιστο μπορεί να είναι μικρότερο από ένα τοπικό ελάχιστο.
        \item Ισχύει ότι $\lim\limits_{x\to0^+}\log(x)=-\infty$.
    \end{alist}
    \item Έστω $f$ παραγωγίσιμη με $f'(x) > 0$ για κάθε $x$ σε ένα διάστημα $\Delta$. Ποιο από τα παρακάτω δεν ισχύει?
    \begin{tasks}[label=\let\textdexiakeraia\relax\alph*.](2)
        \task Η $f$ είναι γνησίως αύξουσα στο $\Delta$.
        \task Η $f$ αντιστρέφεται στο $\Delta$.
        \task Η $f$ είναι κυρτή στο $\Delta$.
        \task Η $f$ δεν παρουσιάζει ακρότατο στο εσωτερικό του $\Delta$.
    \end{tasks}
\end{erwthma}
\end{thema}
%=========== ΤΕΛΟΣ - 11ο ΘΕΜΑ Α ===========


%=========== 12ο ΘΕΜΑ Α ===========
\begin{thema}{A}
\begin{erwthma}
    \item Έστω η συνάρτηση $ f(x)=\sqrt{x},\ x\in[0,+\infty) $. Να αποδείξετε ότι η $ f $ είναι παραγωγίσιμη στο $ (0,+\infty) $ και ισχύει $ f'(x)=\dfrac{1}{2\sqrt{x}} $, δηλαδή $ \left(\sqrt{x}\right)'=\dfrac{1}{2\sqrt{x}} $.
    \item Πότε το σημείο $A(x_0, f(x_0))$ λέγεται σημείο καμπής της γραφικής παράστασης της $f$?
    \item Τι ονομάζουμε τοπικά ακρότατα μιας συνάρτησης $f$?
    \item Να χαρακτηρίσετε τις προτάσεις που ακολουθούν, γράφοντας στο τετράδιό σας, δίπλα στο γράμμα που αντιστοιχεί σε κάθε πρόταση, τη λέξη \textbf{Σωστό}, αν η πρόταση είναι σωστή, ή \textbf{Λάθος}, αν η πρόταση είναι λανθασμένη.
    \begin{alist}
        \item Αν δύο συναρτήσεις $f, g$ έχουν το ίδιο πεδίο ορισμού και ισχύει $f(x)=g(x)$ για κάθε $x$, τότε είναι ίσες.
        \item Αν $\lim\limits_{x\to x_0} f(x) = \ell$, τότε οι τιμές της $f(x)$ προσεγγίζουν το $\ell$ όσο το $x$ πλησιάζει το $x_0$.
        \item Το όριο σύνθεσης $\lim\limits_{x\to x_0} f(g(x))$ ισούται πάντα με $f(\lim\limits_{x\to x_0} g(x))$, αν $\lim\limits_{x\to x_0} g(x)$ υπάρχει και $f$ συνεχής.
        \item Αν $f''(x_0) = 0$, τότε το $x_0$ είναι απαραίτητα σημείο καμπής.
        \item Αν $f, g$ κυρτές σε ένα διάστημα $\varDelta$, το άθροισμά τους $f+g$ είναι κυρτή στο $\varDelta$.
    \end{alist}
    \item Η εξίσωση της εφαπτομένης της $C_f$ στο σημείο $(x_0, f(x_0))$ δίνεται από τη σχέση:
    \begin{tasks}[label=\let\textdexiakeraia\relax\alph*.](2)
        \task $y = f(x_0) + f'(x)(x - x_0)$
        \task $y - f(x_0) = f'(x_0)(x - x_0)$
        \task $y = f'(x_0) \cdot x + f(x_0)$
        \task $y = f(x_0)(x - x_0)$
    \end{tasks}
\end{erwthma}
\end{thema}
%=========== ΤΕΛΟΣ - 12ο ΘΕΜΑ Α ===========


%=========== 13ο ΘΕΜΑ Α ===========
\begin{thema}{A}
\begin{erwthma}
    \item Δίνεται η συνάρτηση $ f(x)=\ln{|x|} $ με πεδίο ορισμού το $ \mathbb{R}^* $. Να δείξετε ότι η $ f $ είναι παραγωγίσιμη στο $ \mathbb{R}^* $ με $ f'(x)=\dfrac{1}{x} $.
    \item Πότε λέμε ότι μια συνάρτηση $f$ στρέφει τα κοίλα προς τα άνω (ή είναι κυρτή) στο $\Delta$?
    \item Πότε μια συνάρτηση $f$ λέγεται παραγωγίσιμη σε ένα ανοικτό διάστημα $(a,\beta)$ και πότε σε ένα κλειστό διάστημα $[a,\beta]$?
    \item Να χαρακτηρίσετε τις προτάσεις που ακολουθούν, γράφοντας στο τετράδιό σας, δίπλα στο γράμμα που αντιστοιχεί σε κάθε πρόταση, τη λέξη \textbf{Σωστό}, αν η πρόταση είναι σωστή, ή \textbf{Λάθος}, αν η πρόταση είναι λανθασμένη.
    \begin{alist}
        \item Αν μια συνάρτηση $f$ ικανοποιεί τις προϋποθέσεις του Θ.Μ.Τ. στο διάστημα $[a,\beta]$ τότε υπάρχει $\xi\in(a,\beta)$ τέτοιο ώστε $f'(\xi)=\frac{f(a)-f(\beta)}{a-\beta}$.
        \item Ισχύει ότι $\dint_{a}^{\beta} (f(x) + g(x)) \d x = \dint_{a}^{\beta} f(x) \d x + \dint_{a}^{\beta} g(x) \d x$, όταν οι $f, g$ είναι συνεχείς στο $[a, \beta]$.
        \item Αν $f'(x_0) = 0$, τότε υποχρεωτικά το $x_0$ είναι τοπικό ακρότατο.
        \item Κάθε συνεχής συνάρτηση στο $\mathbb{R}$ έχει παράγουσα.
        \item Δύο συναρτήσεις είναι ίσες μόνο αν έχουν ίδιους τύπους.
\end{alist}
    \item Έστω $f$ συνεχής στο $[a, \beta]$. Ποια από τις παρακάτω συνθήκες εξασφαλίζει, σύμφωνα με το Θεώρημα \eng{Bolzano}, την ύπαρξη μιας τουλάχιστον ρίζας στο $(a, \beta)$;
    \begin{tasks}[label=\let\textdexiakeraia\relax\alph*.](4)
        \task $f(a) \cdot f(\beta) \ge 0$
        \task $f(a) \cdot f(\beta) < 0$
        \task $f(a) + f(\beta) = 0$
        \task $f(a) - f(\beta) > 0$
    \end{tasks}
\end{erwthma}
\end{thema}
%=========== ΤΕΛΟΣ - 13ο ΘΕΜΑ Α ===========


%=========== 14ο ΘΕΜΑ Α ===========
\begin{thema}{A}
\begin{erwthma}
    \item Να αποδείξετε ότι η συνάρτηση $ f(x)=\ef{x} $ είναι παραγωγίσιμη στο σύνολο $ A=\{x\in\mathbb{R}|\syn{x}\neq 0\} $ και ισχύει $ f'(x)=\dfrac{1}{\syn^2{x}} $.
    \item Τι ορίζουμε ως εφαπτομένη της γραφικής παράστασης μιας συνάρτησης $f$ σε ένα σημείο της $A(x_0, f(x_0))$?
    \item Πώς ορίζεται η αντίστροφη συνάρτηση $f^{-1}$ μιας $1-1$ συνάρτησης $f$?
    \item Να χαρακτηρίσετε τις προτάσεις που ακολουθούν, γράφοντας στο τετράδιό σας, δίπλα στο γράμμα που αντιστοιχεί σε κάθε πρόταση, τη λέξη \textbf{Σωστό}, αν η πρόταση είναι σωστή, ή \textbf{Λάθος}, αν η πρόταση είναι λανθασμένη.
    \begin{alist}
        \item Αν μια παραγωγίσιμη συνάρτηση $f:[a,\beta]\to\mathbb{R}$ είναι $1-1$ τότε δεν ικανοποιεί τις προϋποθέσεις του θεωρήματος \eng{Rolle} στο $[a,\beta]$.
        \item Αν η $f'$ είναι συνεχής συνάρτηση, τότε $\dint_{a}^{\beta} f'(x)g(x)\d x = [f(x)g(x)]_{a}^{\beta} - \dint_{a}^{\beta} f(x)g'(x)\d x$.
        \item Αν $f(x) \ge 0$ στο $[a, \beta]$ και $\dint_a^\beta f(x)\d x = 0$ (με $f$ συνεχής), τότε $f(x) = 0$ σε όλο το $[a, \beta]$.
        \item Αν η $f$ έχει τοπικό ακρότατο στο εσωτερικό σημείο $x_0$ και είναι παραγωγίσιμη σε αυτό, τότε $f'(x_0) = 0$.
        \item Μια πολυωνυμική συνάρτηση βαθμού $\ge 2$ μπορεί να έχει πλάγια ασύμπτωτη.
    \end{alist}
    \item Αν $f'(x_0) = 0$ και $f''(x_0) > 0$, τότε στο $x_0$ η $f$:
    \begin{tasks}[label=\let\textdexiakeraia\relax\alph*.](2)
        \task Παρουσιάζει τοπικό μέγιστο.
        \task Παρουσιάζει τοπικό ελάχιστο.
        \task Παρουσιάζει σημείο καμπής.
        \task Δεν εξάγεται συμπέρασμα.
    \end{tasks}
\end{erwthma}
\end{thema}
%=========== ΤΕΛΟΣ - 14ο ΘΕΜΑ Α ===========


%=========== 15ο ΘΕΜΑ Α ===========
\begin{thema}{A}
\begin{erwthma}
    \item Έστω μια συνάρτηση $ f $ ορισμένη σε ένα διάστημα $ \varDelta $. Αν \begin{itemize} \item $ x_0 $ είναι ένα \textit{εσωτερικό} σημείο του $ \varDelta $ \item η $ f $ παρουσιάζει τοπικό ακρότατο στο $ x_0 $ και \item είναι παραγωγίσιμη στο $ x_0 $ τότε \end{itemize} \[ f'(x_0)=0 \]
    \item Πότε η ευθεία $y=y_0$ ονομάζεται οριζόντια ασύμπτωτη της γραφικής παράστασης μιας συνάρτησης $f$?
    \item Τι ονομάζεται γραφική παράσταση μιας συνάρτησης $f$?
    \item Να χαρακτηρίσετε τις προτάσεις που ακολουθούν, γράφοντας στο τετράδιό σας, δίπλα στο γράμμα που αντιστοιχεί σε κάθε πρόταση, τη λέξη \textbf{Σωστό}, αν η πρόταση είναι σωστή, ή \textbf{Λάθος}, αν η πρόταση είναι λανθασμένη.
    \begin{alist}
        \item Αν η συνάρτηση $f$ έχει κατακόρυφη ασύμπτωτη την ευθεία $x = x_0$, τότε σίγουρα η $f$ δεν είναι συνεχής στο $x_0$.
        \item Σύμφωνα με το Θεώρημα Ενδιάμεσων Τιμών μία συνεχής συνάρτηση $f$ παίρνει όλες τις τιμές μεταξύ $f(a)$ και $f(\beta)$.
        \item Αν η συνάρτηση $f$ είναι γνησίως αύξουσα, τότε υποχρεωτικά $f'(x) > 0$ για κάθε εσωτερικό σημείο $x$.
        \item Το πεδίο ορισμού της $f+g$ είναι το $D_f\cup D_g$.
        \item Για κάθε $x>0$, ισχύει $\ln x^2 = 2\ln x$.
    \end{alist}
    \item Έστω $f(x) \ge 0$. Το εμβαδόν του χωρίου μεταξύ $C_f$, άξονα $x'x$ και ευθειών $x=a, x=\beta$ (με $a<\beta$) δίνεται από:
    \begin{tasks}[label=\let\textdexiakeraia\relax\alph*.](4)
        \task $E = \dint_\beta^a f(x) \d x$
        \task $E = - \dint_a^\beta f(x) \d x$
        \task $E = \dint_a^\beta f(x) \d x$
        \task $E = |f(a) - f(\beta)|$
    \end{tasks}
\end{erwthma}
\end{thema}
%=========== ΤΕΛΟΣ - 15ο ΘΕΜΑ Α ===========


%=========== 16ο ΘΕΜΑ Α ===========
\begin{thema}{A}
\begin{erwthma}
    \item Να διατυπώσετε και να δώσετε τη γεωμετρική ερμηνεία του Θ.Μ.Τ.
    \item Πότε η ευθεία $y=\lambda x+\beta$ ονομάζεται πλάγια ασύμπτωτη της γραφικής παράστασης μιας συνάρτησης $f$ στο $+\infty$?
    \item Τι ονομάζουμε κρίσιμα σημεία μιας συνάρτησης $f$?
    \item Να χαρακτηρίσετε τις προτάσεις που ακολουθούν, γράφοντας στο τετράδιό σας, δίπλα στο γράμμα που αντιστοιχεί σε κάθε πρόταση, τη λέξη \textbf{Σωστό}, αν η πρόταση είναι σωστή, ή \textbf{Λάθος}, αν η πρόταση είναι λανθασμένη.
    \begin{alist}
        \item Αν μια συνάρτηση παίρνει μέγιστη και ελάχιστη τιμή σε ένα διάστημα, τότε είναι συνεχής.
        \item Το πεδίο ορισμού της $f^{-1}$ ταυτίζεται με το σύνολο τιμών της $f$.
        \item Αν $f$ συνεχής και άρτια συνάρτηση με $f(0)<0$ και $\lim\limits_{x\to+\infty}f(x)>0$, τότε έχει τουλάχιστον δύο ρίζες.
        \item $\dint_a^\beta c \d x = c(\beta-a)$ για κάθε σταθερά $c \in \mathbb{R}$.
        \item Αν $0\leq f(x)\leq M$ με $M\in\mathbb{R}$ κοντά στο $x_0$ τότε $\lim\limits_{x\to x_0}{\left[(x-x_0)f(x)\right]}=0$.
    \end{alist}
    \item Αν $f$ συνεχής στο $[a, \beta]$ και $m \le f(x) \le M$ για κάθε $x \in [a, \beta]$, τότε:
    \begin{tasks}[label=\let\textdexiakeraia\relax\alph*.](2)
        \task $\dint_a^\beta f(x) \d x = m(\beta - a)$
        \task $m(\beta-a) \le \dint_a^\beta f(x) \d x \le M(\beta-a)$
        \task $\dint_a^\beta f(x)\d x \ge M(\beta-a)$
        \task Δεν μπορεί να εκτιμηθεί το ολοκλήρωμα.
    \end{tasks}
\end{erwthma}
\end{thema}
%=========== ΤΕΛΟΣ - 16ο ΘΕΜΑ Α ===========


%=========== 17ο ΘΕΜΑ Α ===========
\begin{thema}{A}
\begin{erwthma}
    \item Να αποδείξετε ότι η συνάρτηση $ f(x)=x^{-\nu} $ είναι παραγωγίσιμη στο $ \mathbb{R}^* $ και ισχύει $ f'(x)=-\nu x^{-\nu-1} $, δηλαδή \[ \left(x^{-\nu} \right)'=-\nu x^{-\nu-1} \]
    \item Τι ονομάζουμε αρχική συνάρτηση ή παράγουσα μιας συνάρτησης $f$ σε ένα διάστημα $\Delta$?
    \item Πώς ορίζονται οι πράξεις της διαφοράς $f-g$ και του πηλίκου $\frac{f}{g}$ δύο συναρτήσεων? Ποιο το πεδίο ορισμού τους?
    \item Να χαρακτηρίσετε τις προτάσεις που ακολουθούν, γράφοντας στο τετράδιό σας, δίπλα στο γράμμα που αντιστοιχεί σε κάθε πρόταση, τη λέξη \textbf{Σωστό}, αν η πρόταση είναι σωστή, ή \textbf{Λάθος}, αν η πρόταση είναι λανθασμένη.
    \begin{alist}
        \item Αν $C_f$ έχει οριζόντια ασύμπτωτη στο $+\infty$, δεν μπορεί να έχει και πλάγια ασύμπτωτη στο $+\infty$.
        \item Αν $\dint_a^\beta f(x)\d x > 0$ με $a < \beta$, τότε αναγκαστικά $f(x) > 0$ για κάθε $x \in [a, \beta]$.
        \item Αν $\dint_a^\beta f(x)dx = 0$ και η $f$ δεν είναι παντού μηδέν στο $[a,\beta]$, τότε το εμβαδόν του χωρίου $\Omega$ που περικλείεται από την $C_f$, τον άξονα $x'x$ και τις ευθείες $x=a$ και $x=\beta$ δεν μπορεί να είναι $0$.
        \item Αν $f(x) \le 0$ στο $[a, \beta]$, τότε το εμβαδόν του χωρίου δίνεται από το ολοκλήρωμα $-\dint_a^\beta f(x) \d x$.
        \item Ο συμβολισμός της μεταβλητής ολοκλήρωσης $x$ ή $t$ δεν παίζει ρόλο στην αριθμητική τιμή ενός ορισμένου ολοκληρώματος.
    \end{alist}

    % ----------------- GENERAL THEORY ----------------- %
\end{erwthma}
\end{thema}
%=========== ΤΕΛΟΣ - 17ο ΘΕΜΑ Α ===========


%=========== 18ο ΘΕΜΑ Α ===========
\begin{thema}{A}
\begin{erwthma}
    \item Δίνεται ένα πολυώνυμο $ P(x)=a_\nu x^\nu+a_{\nu-1}x^{\nu-1}+\ldots+a_1x+a_0 $ και $ x_0\in\mathbb{R} $. Να αποδείξετε ότι $ {\displaystyle{\lim_{x\to x_0}{P(x)}=P(x_0)}} $.
    \item Τι ονομάζεται σύνολο τιμών μιας πραγματικής συνάρτησης $f$ με πεδίο ορισμού $A$?
    \item Έστω οι συναρτήσεις $f$ και $g$. Πώς ορίζεται η σύνθεση $g \circ f$ και ποιο είναι το πεδίο ορισμού της?
    \item Να χαρακτηρίσετε τις προτάσεις που ακολουθούν, γράφοντας στο τετράδιό σας, δίπλα στο γράμμα που αντιστοιχεί σε κάθε πρόταση, τη λέξη \textbf{Σωστό}, αν η πρόταση είναι σωστή, ή \textbf{Λάθος}, αν η πρόταση είναι λανθασμένη.
    \begin{alist}
        \item Κάθε συνεχής συνάρτηση με πεδίο ορισμού κλειστό διάστημα $[a, \beta]$ έχει πάντοτε μέγιστη και ελάχιστη τιμή.
        \item Αν η $f$ είναι γνησίως αύξουσα σε ένα διάστημα $\Delta$, τότε η $\dfrac{1}{f}$ είναι σίγουρα γνησίως φθίνουσα στο $\Delta$.
        \item Αν $f'(x) = g'(x)$ τότε $\dint_a^\beta f'(x)\d x = \dint_a^\beta g'(x)\d x \Rightarrow f(\beta)-f(a) = g(\beta)-g(a)$.
        \item Η στιγμιαία ταχύτητα ενός σώματος ισούται με την παράγωγο της θέσης του.
        \item Μια συνεχής συνάρτηση στο $x_0$ ορίζεται απαραίτητα στο $x_0$.
    \end{alist}
    \item Έστω οι συναρτήσεις $f, g$. Το πεδίο ορισμού της σύνθεσης $f \circ g$ ορίζεται ως:
    \begin{tasks}[label=\let\textdexiakeraia\relax\alph*.](2)
        \task $\big\{x \in D_f \mid g(x) \in D_g\big\}$
        \task $\big\{x \in D_g \mid f(x) \in D_f\big\}$
        \task $\big\{x \in D_g \mid g(x) \in D_f\big\}$
        \task $\big\{x \in D_f \mid f(x) \in D_g\big\}$
    \end{tasks}
\end{erwthma}
\end{thema}
%=========== ΤΕΛΟΣ - 18ο ΘΕΜΑ Α ===========


%=========== 19ο ΘΕΜΑ Α ===========
\begin{thema}{A}
\begin{erwthma}
    \item Να αποδείξετε ότι η συνάρτηση $ f(x)=x^{-\nu} $ είναι παραγωγίσιμη στο $ \mathbb{R}^* $ και ισχύει $ f'(x)=-\nu x^{-\nu-1} $, δηλαδή \[ \left(x^{-\nu} \right)'=-\nu x^{-\nu-1} \]
    \item Τι ονομάζουμε πραγματική συνάρτηση με πεδίο ορισμού $A$?
    \item Πότε μια συνάρτηση λέγεται γνησίως φθίνουσα σε ένα διάστημα $\Delta$?
    \item Να χαρακτηρίσετε τις προτάσεις που ακολουθούν, γράφοντας στο τετράδιό σας, δίπλα στο γράμμα που αντιστοιχεί σε κάθε πρόταση, τη λέξη \textbf{Σωστό}, αν η πρόταση είναι σωστή, ή \textbf{Λάθος}, αν η πρόταση είναι λανθασμένη.
    \begin{alist}
        \item Το $\dint_{a}^{\beta} f(x) \d x$ με $a, \beta$ σταθερές, αποτελεί μια συνάρτηση της μεταβλητής $x$.
        \item Το σημείο καμπής είναι το σημείο όπου ο ρυθμός μεταβολής της κλίσης αλλάζει μονοτονία.
        \item Αν $f'(x) > 0$ για κάθε $x \in \Delta$, τότε η $f$ είναι γνησίως αύξουσα στο $\Delta$.
        \item Η εφαπτομένη της $e^x$ στο $0$ είναι η ευθεία $y=x+1$.
        \item Σύμφωνα με το Θ.Μ.Τ., υπάρχει τουλάχιστον ένα $\xi \in (a, \beta)$ ώστε η εφαπτομένη στο $\xi$ να είναι παράλληλη στη χορδή $AB$.
    \end{alist}
    \item To $\displaystyle\lim_{x\to 0} \frac{e^x - 1}{x}$ ισούται με:
    \begin{tasks}[label=\let\textdexiakeraia\relax\alph*.](4)
        \task $0$
        \task $1$
        \task $e$
        \task $+\infty$
    \end{tasks}
\end{erwthma}
\end{thema}
%=========== ΤΕΛΟΣ - 19ο ΘΕΜΑ Α ===========


%=========== 20ο ΘΕΜΑ Α ===========
\begin{thema}{A}
\begin{erwthma}
    \item Έστω η συνάρτηση $f(x)=x,\ x\in\mathbb{R}$. Να αποδείξετε ότι η $f$ είναι παραγωγίσιμη στο $\mathbb{R}$ και ισχύει $f'(x)=1$, δηλαδή $ (x)'=1 $.
    \item Πότε μια συνάρτηση $f:A\to\mathbb{R}$ ονομάζεται συνάρτηση $1-1$?
    \item Τι ονομάζουμε ολικό μέγιστο και τι ολικό ελάχιστο μιας συνάρτησης $f$?
    \item Να χαρακτηρίσετε τις προτάσεις που ακολουθούν, γράφοντας στο τετράδιό σας, δίπλα στο γράμμα που αντιστοιχεί σε κάθε πρόταση, τη λέξη \textbf{Σωστό}, αν η πρόταση είναι σωστή, ή \textbf{Λάθος}, αν η πρόταση είναι λανθασμένη.
    \begin{alist}
        \item Αν $F$ είναι μια παράγουσα της καλώς ορισμένης και συνεχούς $f$, τότε ισχύει $F'(x) = f(x)$ στο πεδίο ορισμού της.
        \item Το ολοκλήρωμα $\dint_1^1 f(x)\d x$ είναι μηδέν για οποιαδήποτε συνεχή συνάρτηση $f$.
        \item Μια συνάρτηση μπορεί να έχει το πολύ δύο οριζόντιες ασύμπτωτες.
        \item Για την απροσδιοριστία $0 \cdot \infty$, μπορούμε να εφαρμόσουμε κανόνα \eng{De L'Hospital}
        \item Αν $f''(x)>0$ στο $\Delta$, τότε το τμήμα της $C_f$ στο $\Delta$ βρίσκεται ολόληρη πάνω από την εφαπτομένη της, με εξαίρεση το σημείο επαφής.
    \end{alist}
    \item Αν $\lim\limits_{x\to x_0} f(x) = \ell\in\mathbb{R}$ και $\lim\limits_{x\to x_0} g(x) = m\in\mathbb{R}$, ποια ιδιότητα δεν ισχύει γενικά?
    \begin{tasks}[label=\let\textdexiakeraia\relax\alph*.](2)
        \task $\lim\limits_{x\to x_0} (f+g) = \ell + m$
        \task $\lim\limits_{x\to x_0} (f \cdot g) = \ell \cdot m$
        \task $\lim\limits_{x\to x_0} (f/g) = \ell/m$
        \task $\lim\limits_{x\to x_0} (f - g) = \ell - m$
    \end{tasks}
\end{erwthma}
\end{thema}
%=========== ΤΕΛΟΣ - 20ο ΘΕΜΑ Α ===========


%=========== 21ο ΘΕΜΑ Α ===========
\begin{thema}{A}
\begin{erwthma}
    \item Αν $ f:A\to\mathbb{R} $ με $ f(x)=\dfrac{P(x)}{Q(x)} $ είναι μια ρητή συνάρτηση και $ x_0\in A $, να αποδείξετε ότι $ {\displaystyle{\lim_{x\to x_0}{\dfrac{P(x)}{Q(x)}}=\dfrac{P(x_0)}{Q(x_0)}}} $ εφόσον $ Q(x_0)\neq0 $.
    \item Πώς ορίζεται η αντίστροφη συνάρτηση $f^{-1}$ μιας $1-1$ συνάρτησης $f$?
    \item Πότε μια συνάρτηση $f$ ονομάζεται συνεχής? Πότε συνεχής σε κλειστό διάστημα $[a, \beta]$?
    \item Να χαρακτηρίσετε τις προτάσεις που ακολουθούν, γράφοντας στο τετράδιό σας, δίπλα στο γράμμα που αντιστοιχεί σε κάθε πρόταση, τη λέξη \textbf{Σωστό}, αν η πρόταση είναι σωστή, ή \textbf{Λάθος}, αν η πρόταση είναι λανθασμένη.
    \begin{alist}
        \item Το Θεώρημα \eng{Rolle} εξασφαλίζει τουλάχιστον μια οριζόντια εφαπτομένη.
        \item Αν η συνάρτηση είναι συνεχής και $1-1$ σε διάστημα, τότε είναι γνησίως μονότονη.
        \item Η γραφική παράσταση της συνάρτησης $-f(x)$ είναι συμμετρική της $C_f$ ως προς τον άξονα $x'x$.
        \item Αν η συνάρτηση $s(t)$ παριστάνει τη θέση κινητού την χρονική στιγμή $t$, τότε ο ρυθμός μεταβολής της ταχύτητας $v(t)$ δίνεται από τον τύπο $a(t)=s''(t)$.
        \item Η μορφή $\dfrac{\pm\infty}{0}$ είναι απροσδιοριστία.
    \end{alist}
    \item Μια ασύμπτωτη ευθεία πλάγιας μορφής ορίζεται από τη σχέση $\lim\limits [\dots] = 0$. Το μέσα όριο είναι:
    \begin{tasks}[label=\let\textdexiakeraia\relax\alph*.](4)
        \task $f(x) \cdot (\lambda x + \mu)$
        \task $f(x) - \lambda x$
        \task $f(x) - \lambda$
        \task $f(x) - (\lambda x + \mu)$
    \end{tasks}
\end{erwthma}
\end{thema}
%=========== ΤΕΛΟΣ - 21ο ΘΕΜΑ Α ===========


%=========== 22ο ΘΕΜΑ Α ===========
\begin{thema}{A}
\begin{erwthma}
    \item Να αποδείξετε ότι η εκθετική συνάρτηση $ f(x)=a^x $ με $ 0<a\neq 1 $ είναι παραγωγίσιμη στο $ \mathbb{R} $ με $ f'(x)=a^x\cdot\ln{a} $.
    \item Τι ονομάζεται γραφική παράσταση μιας συνάρτησης $f$?
    \item Τι ονομάζουμε ρυθμό μεταβολής ενός ποσού $y=f(x)$ ως προς ένα ποσό $x$ σε ένα σημείο $x_0$?
    \item Να χαρακτηρίσετε τις προτάσεις που ακολουθούν, γράφοντας στο τετράδιό σας, δίπλα στο γράμμα που αντιστοιχεί σε κάθε πρόταση, τη λέξη \textbf{Σωστό}, αν η πρόταση είναι σωστή, ή \textbf{Λάθος}, αν η πρόταση είναι λανθασμένη.
    \begin{alist}
        \item Η παράγωγος του αθροίσματος $(f(x)+g(x))'$ ισούται με $f'(x)+g'(x)$.
        \item Το εμβαδόν μεταξύ της συνεχούς $f$ και του $x'x$ από το $a$ στο $\beta$ δίνεται πάντα από τον τύπο $E = \dint_a^\beta f(x) \d x$.
        \item Ο κανόνας \eng{De L'Hospital} εφαρμόζεται μόνο σε μορφές $\dfrac{0}{0}$ ή $\dfrac{\pm \infty}{\pm \infty}$.
        \item Αν μια συνάρτηση $f$ είναι κυρτή σε ένα διάστημα $\varDelta$ και $x_0\in\varDelta$ τότε ισχύει $f(x) \ge f(x_0) + f'(x_0)(x-x_0)$ για κάθε $x\in\varDelta$.
        \item Για κάθε πολυωνυμική συνάρτηση βαθμού μεγαλύτερου του 2, αν έχει τοπικά ακρότατα σε σημεία $x_1<x_2$, τότε θα έχει και σημείο καμπής.
    \end{alist}
    \item Αν $\lim\limits_{x\to x_0} f(x) = -\infty$, τότε για τιμές του $x$ κοντά στο $x_0$, ισχύει:
    \begin{tasks}[label=\let\textdexiakeraia\relax\alph*.](4)
        \task $f(x) > 0$
        \task $f(x) < 0$
        \task $f(x) = 0$
        \task $\lim\limits f(x) \in \mathbb{R}$
    \end{tasks}
\end{erwthma}
\end{thema}
%=========== ΤΕΛΟΣ - 22ο ΘΕΜΑ Α ===========


%=========== 23ο ΘΕΜΑ Α ===========
\begin{thema}{A}
\begin{erwthma}
    \item Έστω μια συνάρτηση $ f $ ορισμένη σε ένα διάστημα $ \varDelta $. Να αποδείξετε ότι αν \begin{itemize} \item η $ f $ είναι συνεχής στο $ \varDelta $ και ισχύει \item $ f'(x)=0 $ σε κάθε \emph{εσωτερικό} σημείο του διαστήματος \end{itemize} τότε η $ f $ είναι σταθερή σε όλο το διάστημα $ \varDelta $.
    \item Πότε δύο συναρτήσεις $f, g$ ονομάζονται ίσες?
    \item Πότε το σημείο $A(x_0, f(x_0))$ λέγεται σημείο καμπής της γραφικής παράστασης της $f$?
    \item Να χαρακτηρίσετε τις προτάσεις που ακολουθούν, γράφοντας στο τετράδιό σας, δίπλα στο γράμμα που αντιστοιχεί σε κάθε πρόταση, τη λέξη \textbf{Σωστό}, αν η πρόταση είναι σωστή, ή \textbf{Λάθος}, αν η πρόταση είναι λανθασμένη.
    \begin{alist}
        \item Δίνονται συναρτήσεις $f,g,h$ για τις οποίες ισχύει $g(x)\leq f(x)\leq h(x)$ για κάθε $x$ σε μια περιοχή γύρο από ένα κοινό σημείο $x_0$. Αν $\lim\limits_{x\to x_0}{g(x)}\neq\lim\limits_{x\to x_0}{h(x)}$ τότε το $\lim\limits_{x\to x_0}{f(x)}$ δεν υπάρχει.
        \item Αν $f'(x_0) = 0$ και $f''(x_0) > 0$, τότε η $f$ έχει τοπικό ελάχιστο στο $x_0$.
        \item Αν για δύο συναρτήσεις $f,g$ ισχύει $f'(x) = g'(x)$ σε ένα διάστημα $\Delta$, τότε $f(x) = g(x) + c$ στο $\Delta$.
        \item Το ορισμένο ολοκλήρωμα μιας συνεχούς συνάρτησης, γεωμετρικά προκύπτει ως το όριο αθροισμάτων εμβαδών απείρων ορθογωνίων.
        \item Ισχύει ότι $\left(ln{2x}\right)'=\left(\ln{3x}\right)'$ για κάθε $x>0$.
    \end{alist}
    \item Μία συνεχής συνάρτηση στο κλειστό $[a, \beta]$:
    \begin{tasks}[label=\let\textdexiakeraia\relax\alph*.](2)
        \task Μπορεί να διαπερνά τον άξονα χωρίς ρίζα.
        \task Είναι $1-1$.
        \task Παίρνει υποχρεωτικά μέγιστη τιμή ($Μ$) και ελάχιστη τιμή ($m$) τιμή.
        \task Δεν έχει τοπικά ακρότατα.
    \end{tasks}
\end{erwthma}
\end{thema}
%=========== ΤΕΛΟΣ - 23ο ΘΕΜΑ Α ===========


%=========== 24ο ΘΕΜΑ Α ===========
\begin{thema}{A}
\begin{erwthma}
    \item Να διατυπώσετε το θεώρημα \eng{Bolzano} και να δώσετε τη γεωμετρική ερμηνεία του.
    \item Πώς ορίζονται οι πράξεις της πρόσθεσης $f+g$ και του πολλαπλασιασμού $f \cdot g$ δύο συναρτήσεων? Ποιο το πεδίο ορισμού τους?
    \item Πότε λέμε ότι μια συνάρτηση $f$ στρέφει τα κοίλα προς τα κάτω (ή είναι κοίλη) στο $\Delta$?
    \item Να χαρακτηρίσετε τις προτάσεις που ακολουθούν, γράφοντας στο τετράδιό σας, δίπλα στο γράμμα που αντιστοιχεί σε κάθε πρόταση, τη λέξη \textbf{Σωστό}, αν η πρόταση είναι σωστή, ή \textbf{Λάθος}, αν η πρόταση είναι λανθασμένη.
    \begin{alist}
        \item Αν το $x_0$ είναι εσωτερικό σημείο ενός διαστήματος $\Delta$ και η $f$ παρουσιάζει τοπικό ακρότατο στο $x_0$, τότε το $x_0$ είναι υποχρεωτικά κρίσιμο σημείο της $f$.
        \item Ισχύει ότι $(f(g(x)))' = f'(g(x)) \cdot g'(x)$.
        \item Αν μια συνάρτηση $f$ είναι αντιστρέψιμη τότε ισχύει $f(f^{-1}(x))=x$ για κάθε $x\in D_f$.
        \item Ισχύει ότι $\dint_{a}^{a} f(x) \d x = 0$, ανεξάρτητα από τον τύπο της συνεχούς συνάρτησης $f$.
        \item Αν $\dint_a^\beta f(x) \d x = \dint_a^\beta g(x) \d x$, με $a \ne \beta$ και $f,g$ συνεχείς, τότε απαραίτητα οι γραφικές τους παραστάσεις έχουν τουλάχιστον ένα κοινό σημείο στο $[a, \beta]$.
    \end{alist}
    \item Δίνεται η γραφική παράσταση της συνάρτησης $f(x)$ με μία οριζόντια ασύμπτωτη:
    \begin{center}
    \begin{tikzpicture}
      \begin{axis}[
          aks_on, belh ar,
          xmin=-2, xmax=6,
          ymin=-1, ymax=4,
          xtick=\empty, ytick={2}, yticklabels={$2$},
          width=8cm, height=5cm
      ]
      \addplot[grafikh parastash, smooth, domain=0.2:6] {2 - 1/x};
      \draw[dashed, line width=0.4mm, red!70!black] (axis cs:-2, 2) -- (axis cs:6, 2);
      \node[fill=white, draw=red!30, rounded corners, inner sep=2pt, text=red!70!black, font=\small, anchor=south west] at (axis cs:4, 2) {$y = 2$};
      \end{axis}
    \end{tikzpicture}
    \end{center}
    Σύμφωνα με το σχήμα, ποιο είναι το όριο $\lim\limits_{x\to +\infty} f(x)$?
    \begin{tasks}[label=\let\textdexiakeraia\relax\alph*.](4)
        \task $-\infty$
        \task $0$
        \task $2$
        \task $+\infty$
    \end{tasks}
\end{erwthma}
\end{thema}
%=========== ΤΕΛΟΣ - 24ο ΘΕΜΑ Α ===========


%=========== 25ο ΘΕΜΑ Α ===========
\begin{thema}{A}
\begin{erwthma}
    \item Δίνεται ένα πολυώνυμο $ P(x)=a_\nu x^\nu+a_{\nu-1}x^{\nu-1}+\ldots+a_1x+a_0 $ και $ x_0\in\mathbb{R} $. Να αποδείξετε ότι $ {\displaystyle{\lim_{x\to x_0}{P(x)}=P(x_0)}} $.
    \item Έστω οι συναρτήσεις $f$ και $g$. Πώς ορίζεται η σύνθεση $g \circ f$ και ποιο είναι το πεδίο ορισμού της?
    \item Πότε η ευθεία $x=x_0$ ονομάζεται κατακόρυφη ασύμπτωτη της γραφικής παράστασης μιας συνάρτησης $f$?
    \item Να χαρακτηρίσετε τις προτάσεις που ακολουθούν, γράφοντας στο τετράδιό σας, δίπλα στο γράμμα που αντιστοιχεί σε κάθε πρόταση, τη λέξη \textbf{Σωστό}, αν η πρόταση είναι σωστή, ή \textbf{Λάθος}, αν η πρόταση είναι λανθασμένη.
    \begin{alist}
        \item Αν η $f$ είναι συνεχής και γνησίως αύξουσα στο $[a, \beta]$, τότε σύνολο τιμών είναι το $[f(a), f(\beta)]$.
        \item Υπάρχει συνεχής συνάρτηση στο $\mathbb{R}$ που δεν τέμνει κανέναν από τους δύο άξονες.
        \item O κανόνας της παραγώγου γινομένου δύο συναρτήσεων είναι $(f(x)g(x))' = f'(x)g'(x)$.
        \item Ισχύει η ιδιότητα: $\dint_a^\beta c f(x) \d x = c \dint_a^\beta f(x) \d x$.
        \item Η συνάρτηση $\dfrac{|x|}{x}$ έχει παράγωγο παντού στο $\mathbb{R}^*$ ίση με μηδέν.
    \end{alist}
    \item Για ποια μορφή \textbf{δεν} μπορούμε να εφαρμόσουμε άμεσα τον κανόνα του \en{L'Hospital}?
    \begin{tasks}[label=\let\textdexiakeraia\relax\alph*.](4)
        \task $\frac{0}{0}$
        \task $\frac{+\infty}{+\infty}$
        \task $\frac{-\infty}{+\infty}$
        \task $0 \cdot \infty$
    \end{tasks}
\end{erwthma}
\end{thema}
%=========== ΤΕΛΟΣ - 25ο ΘΕΜΑ Α ===========


%=========== 26ο ΘΕΜΑ Α ===========
\begin{thema}{A}
\begin{erwthma}
    \item Δίνεται μια συνάρτηση $ f $ η οποία είναι παραγωγίσιμη σ’ ένα διάστημα $ (a,\beta) $, με
εξαίρεση ίσως ένα σημείο του $ x_0 $, στο οποίο όμως είναι συνεχής. Να αποδείξετε ότι αν $ f'(x)<0 $ για κάθε $ x\in(a,x_0) $ και $ f'(x)>0 $ για κάθε $ x\in(x_0,\beta) $ τότε η $ f $ παρουσιάζει τοπικό ελάχιστο στο $ x_0 $.
    \item Πότε μια συνάρτηση λέγεται γνησίως φθίνουσα σε ένα διάστημα $\Delta$?
    \item Πότε η ευθεία $y=y_0$ ονομάζεται οριζόντια ασύμπτωτη της γραφικής παράστασης μιας συνάρτησης $f$?
    \item Να χαρακτηρίσετε τις προτάσεις που ακολουθούν, γράφοντας στο τετράδιό σας, δίπλα στο γράμμα που αντιστοιχεί σε κάθε πρόταση, τη λέξη \textbf{Σωστό}, αν η πρόταση είναι σωστή, ή \textbf{Λάθος}, αν η πρόταση είναι λανθασμένη.
    \begin{alist}
        \item Αν $\lim\limits_{x \to x_0} f(x) = \ell$, τότε η $f$ πρέπει υποχρεωτικά να ορίζεται στο $x_0$.
        \item Κρίσιμα σημεία είναι μόνο όσα μηδενίζουν την παράγωγο.
        \item Η παράγωγος της $\ln x$ είναι $\dfrac{1}{x}$ για $x>0$, και της $\ln |x|$ είναι επίσης $\dfrac{1}{x}$ για $x \ne 0$.
        \item H αντίστροφη μιας γνησίως αύξουσας είναι επίσης γνησίως αύξουσα.
        \item Αν $\dint_a^\beta f(x)\d x > 0$ με $a < \beta$, τότε αναγκαστικά $f(x) > 0$ για κάθε $x \in [a, \beta]$.
    \end{alist}
    \item Για μια συνάρτηση $1-1$ ισχύει απαραίτητα ότι:
    \begin{tasks}[label=\let\textdexiakeraia\relax\alph*.](2)
        \task Είναι γνησίως μονότονη.
        \task $f(x_1) \ne f(x_2) \Rightarrow x_1 \ne x_2$.
        \task $f(x_1) = f(x_2) \Rightarrow x_1 = x_2$.
        \task Tα $C_f, C_{f^{-1}}$ τέμνονται αποκλειστικά στη διχοτόμο $y=x$.
    \end{tasks}
\end{erwthma}
\end{thema}
%=========== ΤΕΛΟΣ - 26ο ΘΕΜΑ Α ===========


%=========== 27ο ΘΕΜΑ Α ===========
\begin{thema}{A}
\begin{erwthma}
    \item Δίνεται μια συνάρτηση $ f $ η οποία είναι συνεχής σε ένα διάστημα $ [a,\beta] $. Να αποδείξετε ότι αν $ G $ είναι μια παράγουσα της $ f $ στο διάστημα $ [a,\beta] $ τότε ισχύει \[ \dint_{a}^{\beta}{f(x)}\d x=G(\beta)-G(a) \]
    \item Τι ονομάζουμε ολικό μέγιστο μιας συνάρτησης $f$?
    \item Πότε η ευθεία $y=\lambda x+\beta$ ονομάζεται πλάγια ασύμπτωτη της γραφικής παράστασης μιας συνάρτησης $f$ στο $+\infty$?
    \item Να χαρακτηρίσετε τις προτάσεις που ακολουθούν, γράφοντας στο τετράδιό σας, δίπλα στο γράμμα που αντιστοιχεί σε κάθε πρόταση, τη λέξη \textbf{Σωστό}, αν η πρόταση είναι σωστή, ή \textbf{Λάθος}, αν η πρόταση είναι λανθασμένη.
    \begin{alist}
        \item Αν $f(x) \ge 2$, τότε $\dint_1^3 f(x) \d x \ge 4$.
        \item Αν μια συνάρτηση είναι συνεχής σε ένα ανοικτό διάστημα, τότε έχει ολικά ακρότατα.
        \item Αν ισχύει $F(x) = G(x) + c$, τότε η $F$ και η $G$ είναι παράγουσες της ίδιας συνάρτησης.
        \item Αν $x_0$ είναι ένα εσωτερικό όπου η $f$ παραγωγίζεται και $f'(x_0) \ne 0$, τότε το $x_0$ δεν είναι ακρότατο.
        \item Η παράγωγος της συνάρτησης $\hm x$ είναι $\syn x$.
    \end{alist}
    \item Αν $f'(x) = 0$ για κάθε $x$ σε ένα \textbf{διάστημα} $\Delta$, τότε:
    \begin{tasks}[label=\let\textdexiakeraia\relax\alph*.](2)
        \task H $f$ είναι γνησίως αύξουσα στο $\Delta$.
        \task H $f$ είναι σταθερή συνάρτηση στο $\Delta$.
        \task H $f$ ταλαντώνεται στο $\Delta$.
        \task H $f$ είναι εξίσωση ευθείας με $a \ne 0$.
    \end{tasks}
\end{erwthma}
\end{thema}
%=========== ΤΕΛΟΣ - 27ο ΘΕΜΑ Α ===========


%=========== 28ο ΘΕΜΑ Α ===========
\begin{thema}{A}
\begin{erwthma}
    \item Έστω η συνάρτηση $ f(x)=x^{\nu},\ \nu\in\mathbb{N}-\{0,1\} $. Να αποδείξετε ότι η $ f $ είναι παραγωγίσιμη στο $ \mathbb{R} $ και ισχύει $ f'(x)=\nu x^{\nu-1} $, δηλαδή $ \left(x^\nu\right)'=\nu x^{\nu-1} $.
    \item Πότε μια συνάρτηση $f$ ονομάζεται συνεχής? Πότε συνεχής σε κλειστό διάστημα $[a, \beta]$?
    \item Τι ονομάζουμε αρχική συνάρτηση ή παράγουσα μιας συνάρτησης $f$ σε ένα διάστημα $\Delta$?
    \item Να χαρακτηρίσετε τις προτάσεις που ακολουθούν, γράφοντας στο τετράδιό σας, δίπλα στο γράμμα που αντιστοιχεί σε κάθε πρόταση, τη λέξη \textbf{Σωστό}, αν η πρόταση είναι σωστή, ή \textbf{Λάθος}, αν η πρόταση είναι λανθασμένη.
    \begin{alist}
        \item Αν $F, G$ είναι δύο παράγουσες της $f$ στο $\Delta$, τότε διαφέρουν κατά μια σταθερά $c$, δηλαδή $F(x) = G(x) + c$.
        \item Οι συναρτήσεις $f(x)=\sqrt[3]{x^4}$ και $g(x)=x^{\frac{4}{3}}$ είναι ίσες.
        \item Οι τριγωνομετρικές συναρτήσεις είναι πάντα συνεχείς στο πεδίο ορισμού τους.
        \item Η ύπαρξη του ορίου $\lim\limits_{x\to x_0} f(x)$ εξαρτάται πάντοτε από την τιμή $f(x_0)$.
        \item Η παράγωγος του $x^x$ είναι $x \cdot x^{x-1}$.
    \end{alist}
    \item Το βασικό τριγωνομετρικό όριο $\lim\limits_{x\to 0} \dfrac{\hm x}{x}$ ισούται με:
    \begin{tasks}[label=\let\textdexiakeraia\relax\alph*.](4)
        \task $0$
        \task $+\infty$
        \task $1$
        \task Δεν υπάρχει
    \end{tasks}
\end{erwthma}
\end{thema}
%=========== ΤΕΛΟΣ - 28ο ΘΕΜΑ Α ===========


%=========== 29ο ΘΕΜΑ Α ===========
\begin{thema}{A}
\begin{erwthma}
    \item Έστω η συνάρτηση $ f(x)=x^{\nu},\ \nu\in\mathbb{N}-\{0,1\} $. Να αποδείξετε ότι η $ f $ είναι παραγωγίσιμη στο $ \mathbb{R} $ και ισχύει $ f'(x)=\nu x^{\nu-1} $, δηλαδή $ \left(x^\nu\right)'=\nu x^{\nu-1} $.
    \item Τι ονομάζουμε ρυθμό μεταβολής ενός ποσού $y=f(x)$ ως προς ένα ποσό $x$ σε ένα σημείο $x_0$?
    \item Πότε μια συνάρτηση $f$ ονομάζεται συνεχής σε ένα σημείο $x_0$ του πεδίου ορισμού της?
    \item Να χαρακτηρίσετε τις προτάσεις που ακολουθούν, γράφοντας στο τετράδιό σας, δίπλα στο γράμμα που αντιστοιχεί σε κάθε πρόταση, τη λέξη \textbf{Σωστό}, αν η πρόταση είναι σωστή, ή \textbf{Λάθος}, αν η πρόταση είναι λανθασμένη.
    \begin{alist}
        \item Αν $f$ συνεχής στο $[a, \beta]$ και παραγωγίσιμη στο $(a, \beta)$ με $f(a)=f(\beta)$, υπάρχει μια τουλάχιστον ρίζα της $f'$ στο $(a, \beta)$.
        \item Κάθε περιττή συνεχής συνάρτηση στο $\mathbb{R}$ έχει ρίζα το $x=0$.
        \item Αν μια συνάρτηση $f$ είναι περιττή και συνεχής, το εμβαδόν του χωρίου μεταξύ της $C_f$, του $x'x$ και των $x=-a, x=a$ ($a>0$) ισούται με $2\dint_0^a |f(x)| \d x$.
        \item Τα κρίσιμα σημεία εντοπίζονται και στα άκρα του πεδίου ορισμού κλειστού διαστήματος.
        \item Αν $\lim\limits_{x\to x_0} f(x) = \ell$, τότε η συνάρτηση ορίζεται υποχρεωτικά στο $x_0$.
    \end{alist}
    \item Αν $f''(x) > 0$ για κάθε $x \in \Delta$, η γραφική παράσταση της $f$ αναφορικά με την εφαπτομένη της:
    \begin{tasks}[label=\let\textdexiakeraia\relax\alph*.](2)
        \task Βρίσκεται χαμηλότερα, με εξαίρεση το σημείο επαφής.
        \task Βρίσκεται ψηλότερα, με εξαίρεση το σημείο επαφής.
        \task Διαπερνάται από αυτή σε κάθε σημείο.
        \task Συμπίπτει με αυτή.
    \end{tasks}
\end{erwthma}
\end{thema}
%=========== ΤΕΛΟΣ - 29ο ΘΕΜΑ Α ===========


%=========== 30ο ΘΕΜΑ Α ===========
\begin{thema}{A}
\begin{erwthma}
    \item Έστω μια συνάρτηση $ f $ η οποία είναι συνεχής σε ένα διάστημα $ \varDelta $. Να αποδείξετε ότι αν ισχύει $ f'(x)<0 $ σε κάθε \textit{εσωτερικό} σημείο του διαστήματος, τότε η $ f $ είναι γνησίως φθίνουσα σε όλο το $ \varDelta $.
    \item Πότε το σημείο $A(x_0, f(x_0))$ λέγεται σημείο καμπής της γραφικής παράστασης της $f$?
    \item Τι ονομάζουμε πραγματική συνάρτηση με πεδίο ορισμού $A$?
    \item Να χαρακτηρίσετε τις προτάσεις που ακολουθούν, γράφοντας στο τετράδιό σας, δίπλα στο γράμμα που αντιστοιχεί σε κάθε πρόταση, τη λέξη \textbf{Σωστό}, αν η πρόταση είναι σωστή, ή \textbf{Λάθος}, αν η πρόταση είναι λανθασμένη.
    \begin{alist}
        \item Αν $\lim\limits_{x \to x_0} f(x) > 0$, τότε $f(x) > 0$ για κάθε $x$ του πεδίου ορισμού της κοντά στο $x_0$.
        \item Το ολοκλήρωμα του αθροίσματος ισούται με το άθροισμα των ολοκληρωμάτων.
        \item Κάθε συνεχής συνάρτηση σε ανοιχτό διάστημα έχει τουλάχιστον ένα κρίσιμο σημείο.
        \item Το $\lim\limits_{x \to x_0} \sqrt{f(x)}$ υπάρχει στο $\mathbb{R}$ αν και μόνο αν $\lim\limits_{x \to x_0} f(x) \ge 0$.
        \item Η ύπαρξη δεύτερης παραγώγου στο $x_0$ εξασφαλίζει ότι το $x_0$ είναι σημείο καμπής, εάν προηγουμένως γνωρίζουμε ότι $f''(x_0)=0$.
    \end{alist}
    \item Το $\dint_a^\beta c \d x$ (όπου $c$ σταθερός πραγματικός αριθμός), ισούται με:
    \begin{tasks}[label=\let\textdexiakeraia\relax\alph*.](4)
        \task $c$
        \task $c\cdot x$
        \task $c(\beta - a)$
        \task $\beta - a$
    \end{tasks}
\end{erwthma}
\end{thema}
%=========== ΤΕΛΟΣ - 30ο ΘΕΜΑ Α ===========


%=========== 31ο ΘΕΜΑ Α ===========
\begin{thema}{A}
\begin{erwthma}
    \item Να αποδείξετε ότι οι γραφικές παραστάσεις $ C_f $ και $ C_{f^{-1}} $ των συναρτήσεων $ f $ και $ f^{-1} $ είναι συμμετρικές ως προς την ευθεία $ y=x $ που διχοτομεί τις γωνίες $ x\hat{O}y $ και $ x'\hat{O}y' $.
    \item Πότε λέμε ότι μια συνάρτηση $f$ στρέφει τα κοίλα προς τα άνω (ή είναι κυρτή) στο $\Delta$?
    \item Πότε μια συνάρτηση $f:A\to\mathbb{R}$ ονομάζεται συνάρτηση $1-1$?
    \item Να χαρακτηρίσετε τις προτάσεις που ακολουθούν, γράφοντας στο τετράδιό σας, δίπλα στο γράμμα που αντιστοιχεί σε κάθε πρόταση, τη λέξη \textbf{Σωστό}, αν η πρόταση είναι σωστή, ή \textbf{Λάθος}, αν η πρόταση είναι λανθασμένη.
    \begin{alist}
        \item Εάν η εφαπτομένη σε ένα σημείο $M$ της γραφικής παράστασης διαπερνά την καμπύλη, τότε το $M$ ονομάζεται σημείο καμπής της.
        \item Το εμβαδόν μεταξύ των γραφικών παραστάσεων δύο συναρτήσεων $f,g$ είναι $\dint_a^\beta |f(x) - g(x)| \d x$.
        \item Ισχύει ότι $\lim\limits_{x \to 0} \dfrac{1}{|x|} = +\infty$.
        \item Υπάρχει συνάρτηση που η παράγωγός της είναι η $f'(x) = |x|$.
        \item Αν η συνάρτηση $f$ είναι συνεχής και δεν μηδενίζεται σε ένα διάστημα $\Delta$, τότε διατηρεί σταθερό πρόσημο στο $\Delta$.
    \end{alist}
    \item Αν $F_1, F_2$ δύο παράγουσες (αρχικές) της ίδιας συνάρτησης $f$ σε ένα διάστημα $\Delta$, τότε:
    \begin{tasks}[label=\let\textdexiakeraia\relax\alph*.](2)
        \task $F_1(x) = F_2(x)$ για κάθε $x \in \Delta$.
        \task $F_1(x) - F_2(x) = c$ (σταθερά) για κάθε $x \in \Delta$.
        \task $F_1'(x) \ne F_2'(x)$ σε κάποιο $x$.
        \task $F_1 \cdot F_2 = c$ (σταθερά).
    \end{tasks}
\end{erwthma}
\end{thema}
%=========== ΤΕΛΟΣ - 31ο ΘΕΜΑ Α ===========


%=========== 32ο ΘΕΜΑ Α ===========
\begin{thema}{A}
\begin{erwthma}
    \item Έστω η ταυτοτική συνάρτηση $ f(x)=x $. Να αποδείξετε ότι η συνάρτηση $ f $ είναι παραγωγίσιμη στο $ \mathbb{R} $ και ισχύει $ f'(x)=1 $, δηλαδή $ (x)'=1 $.
    \item Πότε η ευθεία $x=x_0$ ονομάζεται κατακόρυφη ασύμπτωτη της γραφικής παράστασης μιας συνάρτησης $f$?
    \item Πώς ορίζεται η αντίστροφη συνάρτηση $f^{-1}$ μιας $1-1$ συνάρτησης $f$?
    \item Να χαρακτηρίσετε τις προτάσεις που ακολουθούν, γράφοντας στο τετράδιό σας, δίπλα στο γράμμα που αντιστοιχεί σε κάθε πρόταση, τη λέξη \textbf{Σωστό}, αν η πρόταση είναι σωστή, ή \textbf{Λάθος}, αν η πρόταση είναι λανθασμένη.
    \begin{alist}
        \item Αν $f(x) > 0$ κοντά στο $x_0$, τότε $\lim\limits_{x\to x_0} f(x) > 0$.
        \item H συνάρτηση $e^x$ δεν έχει κατακόρυφη ασύμπτωτη.
        \item Αν $f$ συνεχής και $f(a)f(\beta)=0$, υπάρχει $x_0 \in [a, \beta]$ με $f(x_0)=0$.
        \item Η εφαπτομένη της γραφικής παράστασης μιας συνάρτησης δεν μπορεί να τέμνει την καμπύλη σε άλλα σημεία.
        \item Κάθε σημείο $x_0$ γι το οποίο ισχύει $f'(x_0)=0$ λέγεται κρίσιμο σημείο.
    \end{alist}
    \item Τo ολοκλήρωμα $\dint_1^e \frac{1}{x} \d x$ ισούται με:
    \begin{tasks}[label=\let\textdexiakeraia\relax\alph*.](4)
        \task $e-1$
        \task $1$
        \task $\frac{1}{e}$
        \task $e$
    \end{tasks}
\end{erwthma}
\end{thema}
%=========== ΤΕΛΟΣ - 32ο ΘΕΜΑ Α ===========


%=========== 33ο ΘΕΜΑ Α ===========
\begin{thema}{A}
\begin{erwthma}
    \item Δίνεται μια συνάρτηση $ f $ ορισμένη σε ένα διάστημα $ \varDelta $ και έστω $ F $ μια παράγουσα της $ f $ στο $ \varDelta $. Να αποδείξετε ότι \begin{alist} \item όλες οι συναρτήσεις της μορφής \[ G(x)=F(x)+c \] είναι παράγουσες της $ f $ στο $ \varDelta $ και ότι \item κάθε άλλη παράγουσα $ G $ της $ f $ στο $ \varDelta $ παίρνει τη μορφή \[ G(x)=F(x)+c \] για κάθε $ x\in\varDelta $, με $ c\in\mathbb{R} $. \end{alist}
    \item Πότε η ευθεία $y=y_0$ ονομάζεται οριζόντια ασύμπτωτη της γραφικής παράστασης μιας συνάρτησης $f$?
    \item Τι ονομάζεται γραφική παράσταση μιας συνάρτησης $f$?
    \item Να χαρακτηρίσετε τις προτάσεις που ακολουθούν, γράφοντας στο τετράδιό σας, δίπλα στο γράμμα που αντιστοιχεί σε κάθε πρόταση, τη λέξη \textbf{Σωστό}, αν η πρόταση είναι σωστή, ή \textbf{Λάθος}, αν η πρόταση είναι λανθασμένη.
    \begin{alist}
        \item Αν η $f$ είναι συνεχής στο $[a, \beta]$ και έχει ρίζα, τότε $f(a)f(\beta) < 0$.
        \item Δύο καμπύλες εφάπτονται, αν έχουν κοινό σημείο με κοινή εφαπτομένη δηλαδή $f(x_0)=g(x_0)$ και $f'(x_0)=g'(x_0)$.
        \item Η συνάρτηση $f(x)=|x-1|$ ικανοποιεί τις προϋποθέσεις του θεωρήματος \eng{Rolle} στο διάστημα $[0,2]$.
        \item Ένα σημείο καμπής είναι σημείο όπου η γραφική παράσταση αλλάζει κυρτότητα.
        \item Η τιμή του ολοκληρώματος $\dint_{a}^{\beta} f(x) \d x$ (εφόσον η $f$ είναι συνεχής), είναι πάντοτε ένας θετικός αριθμός.
    \end{alist}
    \item Το σημείο $M(x_0, f(x_0))$ λέγεται σημείο καμπής αν:
    \begin{tasks}[label=\let\textdexiakeraia\relax\alph*.](2)
        \task H $f$ παρουσιάζει τοπικό μέγιστο και τοπικό ελάχιστο.
        \task H $C_f$ έχει εφαπτομένη και αλλάζει κυρτότητα εκατέρωθεν του $x_0$.
        \task $f'(x_0) = 0$ αυστηρά.
        \task $f(x_0) = 0$.
    \end{tasks}
\end{erwthma}
\end{thema}
%=========== ΤΕΛΟΣ - 33ο ΘΕΜΑ Α ===========


%=========== 34ο ΘΕΜΑ Α ===========
\begin{thema}{A}
\begin{erwthma}
    \item Να αποδείξετε ότι η συνάρτηση $ f(x)=x^a $ με $ a\in\mathbb{R}-\mathbb{Z} $ είναι παραγωγίσιμη στο $ (0,+\infty) $ με $ f'(x)=ax^{a-1} $.
    \item Πότε η ευθεία $y=\lambda x+\beta$ ονομάζεται πλάγια ασύμπτωτη της γραφικής παράστασης μιας συνάρτησης $f$ στο $+\infty$?
    \item Πότε δύο συναρτήσεις $f, g$ ονομάζονται ίσες?
    \item Να χαρακτηρίσετε τις προτάσεις που ακολουθούν, γράφοντας στο τετράδιό σας, δίπλα στο γράμμα που αντιστοιχεί σε κάθε πρόταση, τη λέξη \textbf{Σωστό}, αν η πρόταση είναι σωστή, ή \textbf{Λάθος}, αν η πρόταση είναι λανθασμένη.
    \begin{alist}
        \item Η μορφή $0 \cdot (+\infty)$ καταλήγει πάντα στο 0.
        \item Το όριο πολυωνύμου στο $+\infty$ ισούται πάντοτε με το όριο του μεγιστοβάθμιου όρου του.
        \item Αν $\dint_a^\beta f(x) \d x = 0$ και $a \ne \beta$,με $f$ συνεχή, τότε η $f$ έχει τουλάχιστον μια ρίζα στο $[a, \beta]$.
        \item Κάθε εσωτερικό σημείο $x_0$ στο οποίο μηδενίζεται η παράγωγος, δεν είναι υποχρεωτικά θέση τοπικού ακροτάτου.
        \item Αν $F$ είναι μία παράγουσα της παραγωγίσιμης συνάρτησης $f$ και ισχύει $\dint_{a}^{\beta}{F(x)\d x}=0$ τότε η $f$ ικανοποιεί τις προϋποθέσεις του θεωρήματος \eng{Rolle} στο $[a,\beta]$.
    \end{alist}
    \item Ποια είναι η παράγωγος της $f(x) = \ln|g(x)|$, αν $g(x) \ne 0$ και $g$ παραγωγίσιμη?
    \begin{tasks}[label=\let\textdexiakeraia\relax\alph*.](4)
        \task $\frac{1}{g(x)}$
        \task $\frac{g'(x)}{g(x)}$
        \task $\frac{1}{|g(x)|}$
        \task $g'(x) \cdot \ln|g(x)|$
    \end{tasks}
\end{erwthma}
\end{thema}
%=========== ΤΕΛΟΣ - 34ο ΘΕΜΑ Α ===========


%=========== 35ο ΘΕΜΑ Α ===========
\begin{thema}{A}
\begin{erwthma}
    \item Να διατυπώσετε και να δώσετε τη γεωμετρική ερμηνεία του θεωρήματος \eng{Rolle}.
    \item Τι ονομάζουμε αρχική συνάρτηση ή παράγουσα μιας συνάρτησης $f$ σε ένα διάστημα $\Delta$?
    \item Πώς ορίζονται οι πράξεις του γινομένου $f\cdot g$ και του πηλίκου $\dfrac{f}{g}$ δύο συναρτήσεων? Ποιο το πεδίο ορισμού τους?
    \item Να χαρακτηρίσετε τις προτάσεις που ακολουθούν, γράφοντας στο τετράδιό σας, δίπλα στο γράμμα που αντιστοιχεί σε κάθε πρόταση, τη λέξη \textbf{Σωστό}, αν η πρόταση είναι σωστή, ή \textbf{Λάθος}, αν η πρόταση είναι λανθασμένη.
    \begin{alist}
        \item Αν $\lim\limits_{x \to x_0} f(x) = 0$ και η $g$ είναι φραγμένη κοντά στο $x_0$, τότε $\lim\limits_{x \to x_0} [f(x)g(x)] = 0$.
        \item Αν δύο συναρτήσεις $f, g$ έχουν την ίδια παράγωγο στο $\Delta$, τότε οι εφαπτομένες των $C_f, C_g$ στα σημεία με μια κοινή τετμημένη, είναι είτε παράλληλες είτε ταυτίζονται.
        \item Το $\dint_a^\beta f(x) \d x$ είναι ένας πραγματικός αριθμός, που εξαρτάται από τα όρια $a, \beta$ και τη συνάρτηση.
        \item Αν $\lim\limits_{x\to x_0} f(x) > 0$, τότε υποχρεωτικά $f(x) > 0$ για κάθε $x$ κοντά στο $x_0$.
        \item Αν $\lim\limits_{x\to x_0^+} f(x) \ne \lim\limits_{x\to x_0^-} f(x)$, τότε η $f$ δεν είναι συνεχής στο $x_0$.
    \end{alist}
    \item Αν $\dint_a^\beta f(x) \d x = 0$ και η $f$ είναι συνεχής με $f(x) \ge 0$, τότε αναγκαστικά:
    \begin{tasks}[label=\let\textdexiakeraia\relax\alph*.](2)
        \task $f(x) = 0$ για κάθε $x \in [a, \beta]$.
        \task H $f$ είναι σταθερή αλλά όχι μηδενική.
        \task $a = 0$ ή $\beta = 0$.
        \task H $f$ είναι περιττή.
    \end{tasks}
\end{erwthma}
\end{thema}
%=========== ΤΕΛΟΣ - 35ο ΘΕΜΑ Α ===========


%=========== 36ο ΘΕΜΑ Α ===========
\begin{thema}{A}
\begin{erwthma}
    \item Δίνεται μια συνάρτηση $ f $ η οποία είναι παραγωγίσιμη σ’ ένα διάστημα $ (a,\beta) $, με εξαίρεση ίσως ένα σημείο του $ x_0 $, στο οποίο όμως είναι συνεχής. Να αποδείξετε ότι αν $ f'(x)<0 $ για κάθε $ x\in(a,x_0) $ και $ f'(x)>0 $ για κάθε $ x\in(x_0,\beta) $ τότε η $ f $ παρουσιάζει τοπικό ελάχιστο στο $ x_0 $.
    \item Πότε μια συνάρτηση $f$ λέγεται παραγωγίσιμη σε ένα σημείο $x_0$ του πεδίου ορισμού της;
    \item Έστω οι συναρτήσεις $f$ και $g$. Πώς ορίζεται η σύνθεση $g \circ f$ και ποιο είναι το πεδίο ορισμού της?
    \item Να χαρακτηρίσετε τις προτάσεις που ακολουθούν, γράφοντας στο τετράδιό σας, δίπλα στο γράμμα που αντιστοιχεί σε κάθε πρόταση, τη λέξη \textbf{Σωστό}, αν η πρόταση είναι σωστή, ή \textbf{Λάθος}, αν η πρόταση είναι λανθασμένη.
    \begin{alist}
        \item Η κατακόρυφη ασύμπτωτη μιας συνάρτησης μελετάται στο $\pm \infty$.
        \item Αν η εφαπτομένη είναι παράλληλη στον $x'x$, τότε $f'(x_0) = 0$.
        \item Το πεδίο ορισμού της συνάρτησης παράγωγος $f'$ είναι πάντα υποσύνολο του πεδίου ορισμού της $f$.
        \item Αν η εξίσωση $f'(x)=0$ έχει μια μόνο ρίζα σε ένα διάστημα, η $f$ μπορεί να μην παρουσιάζει κανένα ακρότατο στο διάστημα αυτό.
        \item Το όριο $\lim\limits_{x\to 0}{\dfrac{1}{x^{2\nu+1}}}$ υπάρχει στο $\mathbb{R}$ για κάθε $\nu \in \mathbb{N}$.
    \end{alist}
    \item Η εφαπτομένη της γραφικής παράστασης $C_f$ στο $A(x_0, f(x_0))$ τέμνει τον άξονα $x'x$ υπό γωνία $\omega$. Ισχύει πάντοτε:
    \begin{tasks}[label=\let\textdexiakeraia\relax\alph*.](4)
        \task $f'(x_0) = \hm \omega$
        \task $f'(x_0) = \ef \omega$
        \task $f'(x_0) = \omega$
        \task $f'(x_0) = \syn \omega$
    \end{tasks}


    % ----------------- GRAPHICAL (5 Questions) ----------------- %
    \newpage
\end{erwthma}
\end{thema}
%=========== ΤΕΛΟΣ - 36ο ΘΕΜΑ Α ===========


%=========== 37ο ΘΕΜΑ Α ===========
\begin{thema}{A}
\begin{erwthma}
    \item Δίνεται μια συνάρτηση $ f $ η οποία είναι παραγωγίσιμη σ’ ένα διάστημα $ (a,\beta) $, με
εξαίρεση ίσως ένα σημείο του $ x_0 $, στο οποίο όμως είναι συνεχής. Να αποδείξετε ότι αν η $ f' $ διατηρεί το πρόσημό της σε κάθε $ x\in(a,x_0)\cup(x_0,\beta) $ τότε είναι γνησίως μονότονη στο $ (a,\beta) $ και δεν παρουσιάζει τοπικό ακρότατο στο $ x_0 $.
    \item Τι ονομάζουμε πραγματική συνάρτηση με πεδίο ορισμού $A$?
    \item Πότε μια συνάρτηση λέγεται γνησίως φθίνουσα σε ένα διάστημα $\Delta$?
    \item Να χαρακτηρίσετε τις προτάσεις που ακολουθούν, γράφοντας στο τετράδιό σας, δίπλα στο γράμμα που αντιστοιχεί σε κάθε πρόταση, τη λέξη \textbf{Σωστό}, αν η πρόταση είναι σωστή, ή \textbf{Λάθος}, αν η πρόταση είναι λανθασμένη.
    \begin{alist}
        \item Αν η μεταβολή ενός μεγέθους $N$ δίνεται από παραγωγίσιμη συνάρτηση του χρόνου $t$, ο ρυθμός μεταβολής του μεγέθους υπολογίζεται ως η παράγωγος $N'(t)$.
        \item Αν $\lim\limits_{x\to x_0} |f(x)| = 0$, τότε $\lim\limits_{x\to x_0} f(x) = 0$.
        \item Τα εσωτερικά σημεία όπου δεν υπάρχει η παράγωγος, συμπεριλαμβάνονται στα κρίσιμα σημεία.
        \item Η γραφική παράσταση μιας συνάρτησης μπορεί να τέμνει την οριζόντια ασύμπτωτή της.
        \item Η παράγωγος μιας άρτιας συνάρτησης είναι περιττή.
    \end{alist}
    \item Αν η $f$ είναι συνεχής στο $[a, \beta]$ και $\dint_a^\beta f(x)\d x > 0$ (με $a < \beta$), τότε:
    \begin{tasks}[label=\let\textdexiakeraia\relax\alph*.](2)
        \task $f(x) > 0$ σε κάθε $x \in [a, \beta]$.
        \task Υπάρχει $x_0 \in [a, \beta]$ με $f(x_0) > 0$.
        \task $f(a) > 0$ και $f(\beta) > 0$.
        \task H $f$ δεν έχει ρίζες.
    \end{tasks}
\end{erwthma}
\end{thema}
%=========== ΤΕΛΟΣ - 37ο ΘΕΜΑ Α ===========


%=========== 38ο ΘΕΜΑ Α ===========
\begin{thema}{A}
\begin{erwthma}
    \item Να αποδείξετε ότι η συνάρτηση $ f(x)=x^{-\nu} $
είναι παραγωγίσιμη στο $ \mathbb{R}^* $ και ισχύει $ f'(x)=-\nu x^{-\nu-1} $, δηλαδή
\[ \left(x^{-\nu} \right)'=-\nu x^{-\nu-1} \]
    \item Πότε μια συνάρτηση $f:A\to\mathbb{R}$ ονομάζεται συνάρτηση $1-1$?
    \item Τι ονομάζουμε ολικό μέγιστο και τι ολικό ελάχιστο μιας συνάρτησης $f$?
    \item Να χαρακτηρίσετε τις προτάσεις που ακολουθούν, γράφοντας στο τετράδιό σας, δίπλα στο γράμμα που αντιστοιχεί σε κάθε πρόταση, τη λέξη \textbf{Σωστό}, αν η πρόταση είναι σωστή, ή \textbf{Λάθος}, αν η πρόταση είναι λανθασμένη.
    \begin{alist}
        \item Αν στο σημείο καμπής η συνάρτηση παραγωγίζεται, τότε η εφαπτομένη διαπερνά την καμπύλη.
        \item Αν $F, G$ είναι παράγουσες της συνάρτησης $f$ σε σύνολο που αποτελεί ένωση δύο ξένων διαστημάτων, τότε διαφέρουν κατά μία υποχρεωτικά σταθερά.
        \item Η γραφική παράσταση της $f^{-1}$ είναι συμμετρική της $C_f$ ως προς τον άξονα $x'x$.
        \item H συνάρτηση $f(x)=\ln x$ έχει κατακόρυφη ασύμπτωτη τον άξονα $y'y$.
        \item Το ολοκλήρωμα $\dint_1^2 e^x \d x $ ισούται με $ e^2 - e$.
    \end{alist}
    \item Αν $\lim\limits_{x \to 0} \dfrac{f(x)}{x} = 2$, τότε το $\lim\limits_{x \to 0} \dfrac{f(2x)}{x}$ ισούται με:
    \begin{tasks}[label=\let\textdexiakeraia\relax\alph*.](4)
        \task $2$
        \task $4$
        \task $1$
        \task $0$
    \end{tasks}
\end{erwthma}
\end{thema}
%=========== ΤΕΛΟΣ - 38ο ΘΕΜΑ Α ===========


%=========== 39ο ΘΕΜΑ Α ===========
\begin{thema}{A}
\begin{erwthma}
    \item Να αποδείξετε ότι οι γραφικές παραστάσεις $ C_f $ και $ C_{f^{-1}} $ των συναρτήσεων $ f $ και $ f^{-1} $ είναι συμμετρικές ως προς την ευθεία $ y=x $ που διχοτομεί τις γωνίες $ x\hat{O}y $ και $ x'\hat{O}y' $.
    \item Πώς ορίζεται η αντίστροφη συνάρτηση $f^{-1}$ μιας $1-1$ συνάρτησης $f$?
    \item Πότε μια συνάρτηση $f$ ονομάζεται συνεχής? Πότε συνεχής σε κλειστό διάστημα $[a, \beta]$?
    \item Να χαρακτηρίσετε τις προτάσεις που ακολουθούν, γράφοντας στο τετράδιό σας, δίπλα στο γράμμα που αντιστοιχεί σε κάθε πρόταση, τη λέξη \textbf{Σωστό}, αν η πρόταση είναι σωστή, ή \textbf{Λάθος}, αν η πρόταση είναι λανθασμένη.
    \begin{alist}
        \item H $f(x)=|x|$ είναι συνεχής αλλά όχι παραγωγίσιμη στο 0.
        \item Ο συντελεστής διεύθυνσης της εφαπτομένης στο σημείο $(x_0, f(x_0))$ είναι η $f'(x_0)$.
        \item Το πολυώνυμο $P(x)$ έχει πάντοτε όριο στο $x_0$ την τιμή $P(x_0)$.
        \item Σε κάθε τοπικό ακρότατο η παράγωγος πρέπει να μηδενίζεται.
        \item Αν $\dint_{a}^{\beta} f(x) \d x > 0$ με $a < \beta$, τότε η συνεχής συνάρτηση $f$ θα παίρνει απαραιτήτως και θετικές τιμές μέσα στο $[a, \beta]$.
    \end{alist}
    \item Αν $\lim\limits_{x\to x_0} f(x) = 0$ και $\lim\limits_{x\to x_0} g(x) = +\infty$, τότε το $\lim\limits_{x\to x_0} [f(x)g(x)]$ είναι:
    \begin{tasks}[label=\let\textdexiakeraia\relax\alph*.](4)
        \task Πάντα $0$
        \task Πάντα $+\infty$
        \task Πάντα $-\infty$
        \task Μορφή $0 \cdot (\pm \infty)$
    \end{tasks}
\end{erwthma}
\end{thema}
%=========== ΤΕΛΟΣ - 39ο ΘΕΜΑ Α ===========


%=========== 40ο ΘΕΜΑ Α ===========
\begin{thema}{A}
\begin{erwthma}
    \item Να αποδείξετε ότι η συνάρτηση $ f(x)=\syf{x} $ είναι παραγωγίσιμη στο σύνολο $ A=\{x\in\mathbb{R}|\hm{x}\neq 0\} $ και ισχύει $ f'(x)=-\dfrac{1}{\hm^2{x}} $.
    \item Τι ονομάζεται γραφική παράσταση μιας συνάρτησης $f$?
    \item Τι ονομάζουμε ρυθμό μεταβολής ενός ποσού $y=f(x)$ ως προς ένα ποσό $x$ σε ένα σημείο $x_0$?
    \item Να χαρακτηρίσετε τις προτάσεις που ακολουθούν, γράφοντας στο τετράδιό σας, δίπλα στο γράμμα που αντιστοιχεί σε κάθε πρόταση, τη λέξη \textbf{Σωστό}, αν η πρόταση είναι σωστή, ή \textbf{Λάθος}, αν η πρόταση είναι λανθασμένη.
    \begin{alist}
        \item Ισχύει ότι $\dint_{a}^{\beta} f(a) \d x = f(a)(\beta - a)$.
        \item Αν μια συνάρτηση δεν είναι συνεχής στο $x_0$, τότε δεν είναι παραγωγίσιμη σε αυτό.
        \item Το εμβαδόν ενός χωρίου που περικλείεται από την $C_f$, τον άξονα $x'x$ και τις ευθείες $x=a$, $x=\beta$, είναι ένας μη αρνητικός αριθμός.
        \item Κάθε συνεχής συνάρτηση είναι παράγουσα κάποιας άλλης συνάρτησης.
        \item Το θεώρημα \eng{\eng{Bolzano}} ισχύει μόνο για γνήσια μονότονες συναρτήσεις.
    \end{alist}
    \item Στο παρακάτω σχήμα φαίνεται η \textbf{πρώτη παράγωγος} $f'(x)$, η οποία είναι γνησίως φθίνουσα, και τέμνει τον άξονα $x'x$ στο $x=3$.
    \begin{center}
    \begin{tikzpicture}
      \begin{axis}[
          aks_on, belh ar,
          xmin=-1, xmax=5,
          ymin=-2, ymax=4,
          xtick={3}, xticklabels={$3$},
          ytick=\empty,
          width=8cm, height=5cm
      ]
      \path[name path=axis] (axis cs:-1,0) -- (axis cs:5,0);
      \addplot[name path=f, grafikh parastash, domain=-0.5:4.5] {-x+3};
      \node[anchor=south west] at (axis cs:1.4, 2) {$y = f'(x)$};
      \end{axis}
    \end{tikzpicture}
    \end{center}
    Τι συμπεραίνουμε για το ακρότατο της $f$ στη θέση $x=3$?
    \begin{tasks}[label=\let\textdexiakeraia\relax\alph*.](2)
        \task H $f$ παρουσιάζει τοπικό ελάχιστο.
        \task H $f$ παρουσιάζει τοπικό μέγιστο.
        \task Είναι σημείο καμπής.
        \task Δεν είναι ακρότατο αφού η εφαπτομένη της $f'$ δεν είναι μηδέν.
    \end{tasks}
\end{erwthma}
\end{thema}
%=========== ΤΕΛΟΣ - 40ο ΘΕΜΑ Α ===========


%=========== 41ο ΘΕΜΑ Α ===========
\begin{thema}{A}
\begin{erwthma}
    \item Αν οι συναρτήσεις $ f,g $ είναι παραγωγίσιμες στο $ x_0 $, τότε η συνάρτηση $ f+g $ είναι παραγωγίσιμη στο $ x_0 $ και ισχύει: \[ (f+g)'(x_0)=f'(x_0)+g'(x_0) \]
    \item Να διατυπώσετε το κριτήριο παρεμβολής.
    \item Πότε το σημείο $A(x_0, f(x_0))$ λέγεται σημείο καμπής της γραφικής παράστασης της $f$?
    \item Να χαρακτηρίσετε τις προτάσεις που ακολουθούν, γράφοντας στο τετράδιό σας, δίπλα στο γράμμα που αντιστοιχεί σε κάθε πρόταση, τη λέξη \textbf{Σωστό}, αν η πρόταση είναι σωστή, ή \textbf{Λάθος}, αν η πρόταση είναι λανθασμένη.
    \begin{alist}
        \item Αν οι συναρτήσεις $f+g$ και $g$ είναι συνεχείς σε ένα σημείο $x_0$ τότε είναι και η $f$ συνεχής στο $x_0$.
        \item Η παράγωγος της $|x|$ στο $0$ είναι $0$.
        \item Αν μετατοπίσουμε μια θετική συνάρτηση κατακόρυφα κατά $c>0$, το εμβαδόν μεταξύ της καμπύλης, του $x'x$ και των $x=a, x=\beta$ (εφόσον $a<\beta$) αυξάνεται κατά $c(\beta-a)$.
        \item H ασύμπτωτη της $f(x)=\dfrac{1}{x}$ στο $-\infty$ είναι ο άξονας $x'x$.
        \item Αν η συνάρτηση $f$ είναι συνεχής στο $[a, \beta]$ και ισχύει $f(a)f(\beta)<0$, η εξίσωση $f(x)=0$ έχει τουλάχιστον μια ρίζα.
    \end{alist}
    \item Αν $f$ περιττή και ολοκληρώσιμη στο $[-a, a]$ ($a > 0$), τότε $\dint_{-a}^{a} f(x) \d x$ ισούται με:
    \begin{tasks}[label=\let\textdexiakeraia\relax\alph*.](4)
        \task $2 \dint_0^a f(x)\d x$
        \task $0$
        \task $f(a) - f(-a)$
        \task $a$
    \end{tasks}
\end{erwthma}
\end{thema}
%=========== ΤΕΛΟΣ - 41ο ΘΕΜΑ Α ===========


%=========== 42ο ΘΕΜΑ Α ===========
\begin{thema}{A}
\begin{erwthma}
    \item Να αποδείξετε ότι η συνάρτηση $ f(x)=\ef{x} $ είναι παραγωγίσιμη στο σύνολο $ A=\{x\in\mathbb{R}|\syn{x}\neq 0\} $ και ισχύει $ f'(x)=\dfrac{1}{\syn^2{x}} $.
    \item Πώς ορίζονται οι πράξεις της πρόσθεσης $f+g$ και του πολλαπλασιασμού $f \cdot g$ δύο συναρτήσεων? Ποιο το πεδίο ορισμού τους?
    \item Πότε λέμε ότι μια συνάρτηση $f$ στρέφει τα κοίλα προς τα κάτω (ή είναι κοίλη) στο $\Delta$?
    \item Να χαρακτηρίσετε τις προτάσεις που ακολουθούν, γράφοντας στο τετράδιό σας, δίπλα στο γράμμα που αντιστοιχεί σε κάθε πρόταση, τη λέξη \textbf{Σωστό}, αν η πρόταση είναι σωστή, ή \textbf{Λάθος}, αν η πρόταση είναι λανθασμένη.
    \begin{alist}
        \item Ισχύει ότι $\dint_a^\beta f(x) \d x = F(\beta) - F(a)$, όπου $F$ οποιαδήποτε παράγουσα της $f$.
        \item Μια συνάρτηση $f$ από το $A$ στο $\mathbb{R}$ αντιστοιχίζει κάθε στοιχείο του $A$ σε ακριβώς έναν πραγματικό αριθμό.
        \item H $\ef x$ είναι συνεχής σε όλο το $\mathbb{R}$.
        \item Ρυθμός μεταβολής 0 υποδηλώνει ότι το μέγεθος (στιγμιαία) σταματά να μεταβάλλεται.
        \item Αν η παράγωγος μηδενίζεται σε μεμονωμένα σημεία (όχι σε ολόκληρο διάστημα) και παραμένει θετική στα υπόλοιπα, η συνάρτηση είναι γνησίως αύξουσα.
    \end{alist}
    \item Αν $\lim\limits_{x\to +\infty} \frac{f(x)}{x} = \lambda$ και $\lim\limits_{x\to +\infty} (f(x) - \lambda x) = \mu$, η $C_f$ έχει:
    \begin{tasks}[label=\let\textdexiakeraia\relax\alph*.](2)
        \task Κατακόρυφη ασύμπτωτη $x = \lambda$.
        \task Οριζόντια ασύμπτωτη $y = \mu$.
        \task Πλάγια ασύμπτωτη $y = \lambda x + \mu$ στο $+\infty$.
        \task Καμία ασύμπτωτη.
    \end{tasks}
\end{erwthma}
\end{thema}
%=========== ΤΕΛΟΣ - 42ο ΘΕΜΑ Α ===========


%=========== 43ο ΘΕΜΑ Α ===========
\begin{thema}{A}
\begin{erwthma}
    \item Δίνεται μια συνάρτηση $ f $ η οποία είναι συνεχής σε ένα διάστημα $ [a,\beta] $. Να αποδείξετε ότι αν $ G $ είναι μια παράγουσα της $ f $ στο διάστημα $ [a,\beta] $ τότε ισχύει \[ \dint_{a}^{\beta}{f(x)}\d x=G(\beta)-G(a) \]
    \item Έστω οι συναρτήσεις $f$ και $g$. Πώς ορίζεται η σύνθεση $g \circ f$ και ποιο είναι το πεδίο ορισμού της?
    \item Πότε η ευθεία $x=x_0$ ονομάζεται κατακόρυφη ασύμπτωτη της γραφικής παράστασης μιας συνάρτησης $f$?
    \item Να χαρακτηρίσετε τις προτάσεις που ακολουθούν, γράφοντας στο τετράδιό σας, δίπλα στο γράμμα που αντιστοιχεί σε κάθε πρόταση, τη λέξη \textbf{Σωστό}, αν η πρόταση είναι σωστή, ή \textbf{Λάθος}, αν η πρόταση είναι λανθασμένη.
    \begin{alist}
        \item Αν $\lim\limits_{x\to x_0} f(x) = 0$ και $\lim\limits_{x\to x_0} g(x) = +\infty$, το $\lim\limits_{x\to x_0} (f(x)g(x))$ είναι πάντα 0.
        \item Το Θ.Μ.Τ. είναι ειδική περίπτωση του θεωρήματος \eng{Rolle}.
        \item Αν $f(x) \geq g(x)$ για $x \in [a, \beta]$, τότε το εμβαδόν ανάμεσά τους στο διάστημα $[a,\beta]$ ισούται με $\dint_{a}^{\beta} (f(x)-g(x)) \d x$.
        \item Σε κάθε διάστημα που η γραφική παράσταση μιας συνάρτησης $f$ είναι κοίλη, η παράγωγος $f'$ είναι γνησίως φθίνουσα.
        \item Υπάρχει συνάρτηση που είναι συνεχής στο $\mathbb{R}$ αλλά παρουσιάζει σημείο όπου η εφαπτομένη της δεν ορίζεται.
    \end{alist}
    \item Το ολοκλήρωμα $\dint_{a}^{\beta} f(\hm(x)) \syn(x) \d x$ αντιμετωπίζεται κάνοντας:
    \begin{tasks}[label=\let\textdexiakeraia\relax\alph*.](2)
        \task Παραγοντική ολοκλήρωση.
        \task Αντικατάσταση $u = \hm x$.
        \task Αντικατάσταση $u = \syn x$.
        \task Αδύνατον.
    \end{tasks}
\end{erwthma}
\end{thema}
%=========== ΤΕΛΟΣ - 43ο ΘΕΜΑ Α ===========


%=========== 44ο ΘΕΜΑ Α ===========
\begin{thema}{A}
\begin{erwthma}
    \item {Έστω η συνάρτηση $ f(x)=\sqrt{x},\ x\in[0,+\infty) $. Να αποδείξετε ότι η $ f $ είναι παραγωγίσιμη στο $ (0,+\infty) $ και ισχύει $ f'(x)=\dfrac{1}{2\sqrt{x}} $, δηλαδή $ \left(\sqrt{x}\right)'=\dfrac{1}{2\sqrt{x}} $.}
    \item Πότε μια συνάρτηση λέγεται γνησίως φθίνουσα σε ένα διάστημα $\Delta$?
    \item Πότε η ευθεία $y=y_0$ ονομάζεται οριζόντια ασύμπτωτη της γραφικής παράστασης μιας συνάρτησης $f$?
    \item Να χαρακτηρίσετε τις προτάσεις που ακολουθούν, γράφοντας στο τετράδιό σας, δίπλα στο γράμμα που αντιστοιχεί σε κάθε πρόταση, τη λέξη \textbf{Σωστό}, αν η πρόταση είναι σωστή, ή \textbf{Λάθος}, αν η πρόταση είναι λανθασμένη.
    \begin{alist}
        \item Κάθε συνεχής συνάρτηση σε ένα κλειστό διάστημα $[a, \beta]$ είναι ολοκληρώσιμη σε αυτό.
        \item Μια συνάρτηση $1-1$, της οποίας το πεδίο ορισμού είναι διάστημα, είναι απαραίτητα συνεχής.
        \item Το θεώρημα {\eng{Fermat}} αναφέρεται αποκλειστικά σε ολικά ακρότατα.
        \item Αν $f(x)<g(x)$ για κάθε $x \in [a, \beta]$ (οπού $f,g$ συνεχείς), τότε το εμβαδόν μεταξύ των $C_f, C_g$ δίνεται από το $E = \dint_a^\beta (g(x)-f(x))\d x$, ανεξάρτητα από το αν οι $f, g$ λαμβάνουν αρνητικές τιμές.
        \item Μια συνάρτηση μπορεί να έχει μέγιστο σε σημείο μη παραγωγισιμότητας.
    \end{alist}
    \item Στο παρακάτω σχήμα απεικονίζεται η \textbf{δεύτερη παράγωγος} $f''(x)$. 
    \begin{center}
    \begin{tikzpicture}
      \begin{axis}[
          aks_on, belh ar,
          xmin=-2, xmax=4,
          ymin=-2, ymax=3,
          xtick={1}, xticklabels={$1$},
          ytick=\empty,
          width=8cm, height=5cm
      ]
      \path[name path=axis] (axis cs:-2,0) -- (axis cs:4,0);
      \addplot[name path=f, grafikh parastash, domain=-1:3.5, samples=2] {x-1};
      \node[anchor=north west] at (axis cs:1.5, 2) {$y = f''(x)$};
      \end{axis}
    \end{tikzpicture}
    \end{center}
    Ποιος ισχυρισμός για την αρχική συνάρτηση $f(x)$ είναι σωστός (εφόσον ορίζεται η εφαπτομένη της στο 1)?
    \begin{tasks}[label=\let\textdexiakeraia\relax\alph*.](2)
        \task Η $f$ παρουσιάζει τοπικό ελάχιστο στο $x_0 = 1$.
        \task Το σημείο με τετμημένη $x_0 = 1$ είναι θέση σημείου καμπής της $f$.
        \task Η $f$ είναι κοίλη (στρέφει τα κοίλα κάτω) στο $[1, +\infty)$.
        \task Η πρώτη παράγωγος είναι αδύνατον να βρεθεί.
    \end{tasks}
\end{erwthma}
\end{thema}
%=========== ΤΕΛΟΣ - 44ο ΘΕΜΑ Α ===========


%=========== 45ο ΘΕΜΑ Α ===========
\begin{thema}{A}
\begin{erwthma}
    \item Αν $ f:A\to\mathbb{R} $ με $ f(x)=\dfrac{P(x)}{Q(x)} $ είναι μια ρητή συνάρτηση και $ x_0\in A $, να αποδείξετε ότι $ {\displaystyle{\lim_{x\to x_0}{\dfrac{P(x)}{Q(x)}}=\dfrac{P(x_0)}{Q(x_0)}}} $ εφόσον $ Q(x_0)\neq0 $.
    \item Τι ονομάζουμε ολικό μέγιστο και τι ολικό ελάχιστο μιας συνάρτησης $f$?
    \item Πότε η ευθεία $y=\lambda x+\beta$ ονομάζεται πλάγια ασύμπτωτη της γραφικής παράστασης μιας συνάρτησης $f$ στο $+\infty$?
    \item Να χαρακτηρίσετε τις προτάσεις που ακολουθούν, γράφοντας στο τετράδιό σας, δίπλα στο γράμμα που αντιστοιχεί σε κάθε πρόταση, τη λέξη \textbf{Σωστό}, αν η πρόταση είναι σωστή, ή \textbf{Λάθος}, αν η πρόταση είναι λανθασμένη.
    \begin{alist}
        \item Η παράγωγος της $e^x$ είναι $e^x$, ενώ του $a^x$ είναι $a^x \ln a$.
        \item Το όριο της $\dfrac{1}{x^2}$ στο 0 είναι $+\infty$.
        \item Αν $\lim\limits_{x\to x_0} f(x) = \ell$ και $\lim\limits_{x\to x_0} g(x) = +\infty$, τότε $\lim\limits_{x\to x_0} (f(x)+g(x)) = +\infty$.
        \item Ισχύει ότι $\lim\limits_{x \to +\infty} \dfrac{\hm x}{x} = 1$.
        \item Αν η $f$ είναι $1-1$, τότε κάθε οριζόντια ευθεία τέμνει την $C_f$ το πολύ σε ένα σημείο.
    \end{alist}
    \item Αν $f$ συνεχής στο $[a,\beta]$ με $f(a) \cdot f(\beta) < 0$, τότε:
    \begin{tasks}[label=\let\textdexiakeraia\relax\alph*.](2)
        \task Η $f$ είναι μονότονη.
        \task Η εξίσωση $f(x) = 0$ έχει τουλάχιστον μια ρίζα στο $(a,\beta)$.
        \task Η $f$ μηδενίζεται ακριβώς μία φορά.
        \task Η $f$ δεν είναι παραγωγίσιμη στο $(a,\beta)$.
    \end{tasks}
\end{erwthma}
\end{thema}
%=========== ΤΕΛΟΣ - 45ο ΘΕΜΑ Α ===========


%=========== 46ο ΘΕΜΑ Α ===========
\begin{thema}{A}
\begin{erwthma}
    \item Έστω μια συνάρτηση $ f $ η οποία είναι συνεχής σε ένα διάστημα $ \varDelta $. Να αποδείξετε ότι αν ισχύει $ f'(x)>0 $ σε κάθε \textit{εσωτερικό} σημείο του διαστήματος, τότε η $ f $ είναι γνησίως αύξουσα σε όλο το $ \varDelta $.
    \item Πότε μια συνάρτηση $f$ ονομάζεται συνεχής σε ένα σημείο $x_0$ του πεδίου ορισμού της?
    \item Τι ονομάζουμε αρχική συνάρτηση ή παράγουσα μιας συνάρτησης $f$ σε ένα διάστημα $\Delta$?
    \item Να χαρακτηρίσετε τις προτάσεις που ακολουθούν, γράφοντας στο τετράδιό σας, δίπλα στο γράμμα που αντιστοιχεί σε κάθε πρόταση, τη λέξη \textbf{Σωστό}, αν η πρόταση είναι σωστή, ή \textbf{Λάθος}, αν η πρόταση είναι λανθασμένη.
    \begin{alist}
        \item Κάθε συνάρτηση $1-1$ είναι υποχρεωτικά γνησίως μονότονη.
        \item Αν $f'(x) = 0$ σε ένα διάστημα $\Delta$, τότε η $f$ είναι υποχρεωτικά σταθερή στο $\Delta$.
        \item Αν το $\lim\limits_{x\to +\infty} (f(x) - (x+2)) = 0$, τότε η ευθεία $y = x+2$ είναι πλάγια ασύμπτωτη της $C_f$ στο $+\infty$.
        \item Ισχύει γενικά ότι αν $\lim\limits_{x\to x_0} (f(x)g(x)) = 0$, τότε $\lim\limits_{x\to x_0} f(x) = 0$ ή $\lim\limits_{x\to x_0} g(x) = 0$.
        \item Η συνάρτηση $f(x) = x^3$ είναι γνησίως αύξουσα στο $\mathbb{R}$, παρόλο που η παράγωγός της μηδενίζεται στο $x=0$.
    \end{alist}
    \item Αν $f$ γνησίως αύξουσα και συνεχής στο $[a, \beta]$, τότε το σύνολο τιμών της $f$ είναι:
    \begin{tasks}[label=\let\textdexiakeraia\relax\alph*.](4)
        \task $(f(a), f(\beta))$
        \task $[f(a), f(\beta)]$
        \task $\{f(a), f(\beta)\}$
        \task $\mathbb{R}$
    \end{tasks}
\end{erwthma}
\end{thema}
%=========== ΤΕΛΟΣ - 46ο ΘΕΜΑ Α ===========


%=========== 47ο ΘΕΜΑ Α ===========
\begin{thema}{A}
\begin{erwthma}
    \item Να αποδείξετε ότι αν μια συνάρτηση $f$ είναι παραγωγίσιμη σ’ ένα σημείο \bmath{$ x_0 $} τότε είναι και συνεχής στο σημείο αυτό.
    \item Τι ονομάζουμε ρυθμό μεταβολής ενός ποσού $y=f(x)$ ως προς ένα ποσό $x$ σε ένα σημείο $x_0$?
    \item Ποιες είναι οι πιθανές θέσεις τοπικών ακροτάτων μια συνάρτησης $f$ σε ένα διάστημα $\Delta$? Ποια σημεία ονομάζονται κρίσιμα σημεία?
    \item Να χαρακτηρίσετε τις προτάσεις που ακολουθούν, γράφοντας στο τετράδιό σας, δίπλα στο γράμμα που αντιστοιχεί σε κάθε πρόταση, τη λέξη \textbf{Σωστό}, αν η πρόταση είναι σωστή, ή \textbf{Λάθος}, αν η πρόταση είναι λανθασμένη.
    \begin{alist}
        \item Για τον υπολογισμό του εμβαδού του χωρίου μεταξύ των $C_f, C_g$, είναι απαραίτητο να εντοπίσουμε πρώτα τις τετμημένες των σημείων τομής τους.
        \item Αν $\dint_a^\beta f(x)\d x = 0$, τότε αναγκαστικά $a = \beta$.
        \item Αν $f''(x) > 0$, η $f'$ είναι $1-1$.
        \item Αν $\lim\limits_{x\to x_0} f(x) = +\infty$, τότε $f(x) > 0$ κοντά στο $x_0$.
        \item Το ολοκλήρωμα $\dint_{a}^{\beta} f(x) \d x$ εκφράζει υποχρεωτικά εμβαδόν κάποιου χωρίου.
    \end{alist}
    \item Σύμφωνα με το θεώρημα ενδιάμεσων τιμών, η $f(x)$ παίρνει:
    \begin{tasks}[label=\let\textdexiakeraia\relax\alph*.](2)
        \task Μόνο τιμές γύρω από το μηδέν.
        \task Όλες τις τιμές μεταξύ $f(a)$ και $f(\beta)$ (ως συνεχής).
        \task Μέγιστη και ελάχιστη τιμή.
        \task Μόνο θετικές τιμές.
    \end{tasks}
\end{erwthma}
\end{thema}
%=========== ΤΕΛΟΣ - 47ο ΘΕΜΑ Α ===========


%=========== 48ο ΘΕΜΑ Α ===========
\begin{thema}{A}
\begin{erwthma}
    \item Έστω μια συνάρτηση $ f $ ορισμένη σε ένα διάστημα $ \varDelta $. Αν \begin{itemize} \item $ x_0 $ είναι ένα \textit{εσωτερικό} σημείο του $ \varDelta $ \item η $ f $ παρουσιάζει τοπικό ακρότατο στο $ x_0 $ και \item είναι παραγωγίσιμη στο $ x_0 $ τότε \end{itemize} \[ f'(x_0)=0 \]
    \item Πότε το σημείο $A(x_0, f(x_0))$ λέγεται σημείο καμπής της γραφικής παράστασης της $f$?
    \item Τι ονομάζουμε πραγματική συνάρτηση με πεδίο ορισμού $A$?
    \item Να χαρακτηρίσετε τις προτάσεις που ακολουθούν, γράφοντας στο τετράδιό σας, δίπλα στο γράμμα που αντιστοιχεί σε κάθε πρόταση, τη λέξη \textbf{Σωστό}, αν η πρόταση είναι σωστή, ή \textbf{Λάθος}, αν η πρόταση είναι λανθασμένη.
    \begin{alist}
        \item Αν μια συνάρτηση ορίζεται στο $x_0$ και έχει όριο στο $x_0$, τότε είναι συνεχής στο $x_0$.
        \item Η μορφή $+\infty - \infty$ είναι απροσδιόριστη.
        \item Οι συναρτήσεις $f\circ g$ και $g\circ f$ είναι πάντα ίσες.
        \item Αν $f(x)<g(x)$ κοντά στο $x_0$, μπορεί στο όριο να ισχύει $\lim\limits_{x\to x_0} f(x) = \lim\limits_{x\to x_0} g(x)$.
        \item Αν $f(x) \ge 0$ για $x \in [a, \beta]$, τότε $\dint_a^\beta f(x) \d x \ge 0$.
    \end{alist}
    \item Ποια συνθήκη δεν ανήκει στο Θεώρημα του \eng{Rolle} για το διάστημα $[a, \beta]$?
    \begin{tasks}[label=\let\textdexiakeraia\relax\alph*.](2)
        \task H $f$ είναι συνεχής στο $[a, \beta]$.
        \task H $f$ είναι παραγωγίσιμη στο $[a, \beta]$.
        \task $f(a) = f(\beta)$.
        \task H $f$ είναι παραγωγίσιμη στο $(a, \beta)$.
    \end{tasks}
\end{erwthma}
\end{thema}
%=========== ΤΕΛΟΣ - 48ο ΘΕΜΑ Α ===========


%=========== 49ο ΘΕΜΑ Α ===========
\begin{thema}{A}
\begin{erwthma}
    \item Να αποδείξετε ότι η εκθετική συνάρτηση $ f(x)=a^x $ με $ 0<a\neq 1 $ είναι παραγωγίσιμη στο $ \mathbb{R} $ με $ f'(x)=a^x\cdot\ln{a} $.
    \item Πότε λέμε ότι μια συνάρτηση $f$ στρέφει τα κοίλα προς τα άνω (ή είναι κυρτή) στο $\Delta$?
    \item Πότε μια συνάρτηση $f:A\to\mathbb{R}$ ονομάζεται συνάρτηση $1-1$?
    \item Να χαρακτηρίσετε τις προτάσεις που ακολουθούν, γράφοντας στο τετράδιό σας, δίπλα στο γράμμα που αντιστοιχεί σε κάθε πρόταση, τη λέξη \textbf{Σωστό}, αν η πρόταση είναι σωστή, ή \textbf{Λάθος}, αν η πρόταση είναι λανθασμένη.
    \begin{alist}
        \item Το πεδίο ορισμού της σύνθεσης $f \circ g$ ταυτίζεται πάντοτε με την τομή των πεδίων ορισμού της $f$ και της $g$.
        \item Η κατακόρυφη ευθεία $x=x_0$ λέγεται κατακόρυφη ασύμπτωτη αν έστω ένα πλευρικό όριο στο $x_0$ είναι $+\infty$ ή $-\infty$.
        \item Η $\hm x$ είναι κυρτή στο $(0, \pi)$.
        \item Το πεδίο ορισμού της $f^{-1}$ είναι το σύνολο τιμών της $f$ και το σύνολο τιμών της $f^{-1}$ είναι το πεδίο ορισμού της $f$.
        \item Αν $\dint_a^\beta f(x)\d x = 0$, τότε αναγκαστικά $a = \beta$.
    \end{alist}
    \item Παρατηρήστε την παρακάτω γραφική παράσταση της \textbf{παραγώγου} $f'(x)$ μιας συνεχούς συνάρτησης $f$.
    \begin{center}
    \begin{tikzpicture}
      \begin{axis}[
          aks_on, belh ar,
          xmin=-3, xmax=4,
          ymin=-3, ymax=5,
          xtick={-1, 2}, xticklabels={$-1$, $2$},
          ytick=\empty,
          width=8cm, height=5cm
      ]
      \path[name path=axis] (axis cs:-3,0) -- (axis cs:4,0);
      \addplot[name path=f, grafikh parastash, smooth, domain=-2:3] {(x+1)*(x-2)};
      \node[anchor=south west] at (axis cs:1.5, 1) {$y = f'(x)$};
      \end{axis}
    \end{tikzpicture}
    \end{center}
    Σε ποιο από τα παρακάτω διαστήματα η αρχική συνάρτηση $f(x)$ είναι γνησίως φθίνουσα?
    \begin{tasks}[label=\let\textdexiakeraia\relax\alph*.](4)
        \task $[ -1, 2 ]$
        \task $(-\infty, -1]$
        \task $[2, +\infty)$
        \task $\mathbb{R}$
    \end{tasks}
\end{erwthma}
\end{thema}
%=========== ΤΕΛΟΣ - 49ο ΘΕΜΑ Α ===========


%=========== 50ο ΘΕΜΑ Α ===========
\begin{thema}{A}
\begin{erwthma}
    \item Δίνονται δύο συναρτήσεις $ f,g $ ορισμένες σε ένα διάστημα $ \varDelta $. Να αποδείξετε ότι αν
\begin{itemize}
\item οι συναρτήσεις $ f,g $ είναι συνεχείς στο διάστημα $ \varDelta $ και
\item $ f'(x)=g'(x) $ σε κάθε \emph{εσωτερικό} σημείο του $ \varDelta $
\end{itemize}
τότε υπάρχει σταθερά $ c $ τέτοια ώστε να ισχύει $ f(x)=g(x)+c $ για κάθε $ x\in\varDelta $.
    \item Πότε η ευθεία $x=x_0$ ονομάζεται κατακόρυφη ασύμπτωτη της γραφικής παράστασης μιας συνάρτησης $f$?
    \item Πώς ορίζεται η αντίστροφη συνάρτηση $f^{-1}$ μιας $1-1$ συνάρτησης $f$?
    \item Να χαρακτηρίσετε τις προτάσεις που ακολουθούν, γράφοντας στο τετράδιό σας, δίπλα στο γράμμα που αντιστοιχεί σε κάθε πρόταση, τη λέξη \textbf{Σωστό}, αν η πρόταση είναι σωστή, ή \textbf{Λάθος}, αν η πρόταση είναι λανθασμένη.
    \begin{alist}
        \item Αν η συνάρτηση $f$ είναι συνεχής στο κλειστό διάστημα $[a, \beta]$, παραγωγίσιμη στο $(a, \beta)$ και ισχύει $f(a) \neq f(\beta)$, τότε δεν υπάρχει $\xi \in (a, \beta)$ τέτοιο ώστε $f'(\xi) = 0$.
\item Αν μια συνάρτηση είναι $1-1$, τότε οι γραφικές παραστάσεις $C_f$ και $C_{f^{-1}}$ έχουν κοινά σημεία μόνο πάνω στην ευθεία $y=x$.
\item Έστω συνάρτηση $f$ ορισμένη σε ένα διάστημα $\Delta$ και $x_0$ ένα εσωτερικό σημείο του $\Delta$. Αν η $f$ είναι παραγωγίσιμη στο $x_0$ και το $x_0$ δεν είναι κρίσιμο σημείο, τότε η $f$ δεν παρουσιάζει τοπικό ακρότατο στο $x_0$.
\item Ισχύει ότι $\dint_{a}^{\beta} f(x)g'(x)dx = f(\beta)g(\beta) - f(a)g(a) - \dint_{a}^{\beta} f'(x)g(x)dx$ εφόσον οι $f', g'$ είναι συνεχείς.
\item Αν $\lim\limits_{x \to x_0} f(x) = +\infty$, τότε η $f$ παίρνει μόνο θετικές τιμές σε μια περιοχή του $x_0$.
    \end{alist}
    \item Δίνεται το εμβαδόν ενός κλειστού χωρίου $\Omega$ ανάμεσα στην καμπύλη της $f$, τον $x'x$ και τις ευθείες $x=0, x=4$.
    \begin{center}
    \begin{tikzpicture}
      \begin{axis}[
          aks_on, belh ar,
          xmin=-1, xmax=5,
          ymin=-1, ymax=3,
          xtick={0, 4}, xticklabels={$0$, $4$},
          ytick=\empty,
          width=8cm, height=5cm
      ]
      \addplot[name path=f, grafikh parastash, smooth, domain=0:4] {0.5*x*(4-x)};
      \path[name path=axis] (axis cs:0,0) -- (axis cs:4,0);
      \addplot[maincolor!30, opacity=0.6] fill between[of=f and axis];
      \node[fill=white, draw=maincolor!50, rounded corners, inner sep=3pt, text=maincolor!80!black, font=\small\bfseries] at (axis cs:2, 0.7) {$\Omega = 10$};
      \end{axis}
    \end{tikzpicture}
    \end{center}
    Βασιζόμενοι στο σχήμα (όπου $f(x)>0$ στο $(0,4)$), αν $F$ μια αρχική της $f$, ισχύει:
    \begin{tasks}[label=\let\textdexiakeraia\relax\alph*.](2)
        \task $F(4) - F(0) = -10$
        \task $F'(4) - F'(0) = 10$
        \task $\dint_0^4 f(x) \d x = 10 \Rightarrow F(4) - F(0) = 10$
        \task $\dint_4^0 f(x) \d x = 10$
    \end{tasks}
\end{erwthma}
\end{thema}
%=========== ΤΕΛΟΣ - 50ο ΘΕΜΑ Α ===========